%% !TEX root = ../manual_gaolv5.tex
% Chapter 15.
%% Chapter 15: Moving a program from GAOL 4 to GAOL v5. Sources: gaol.tex of f5bf180
%% lines 6512-6639 (the chapter of the same name), doc/differences.md (whole).
%%
%% Copyright (c) 2026 ENSTA, France
%%
%% By: Jordan NININ (created 2026-11-06)
%%
%% GAOL is a software distributed WITHOUT ANY WARRANTY. Read the associated
%% COPYING.LIB file for information.
\chapter{Moving a program from GAOL~4 to GAOL v5}\label{chap:changes}
\cindex{GAOL~4!moving from}%
\cindex{migration from GAOL~4}%

GAOL v5 continues the GAOL 4.2.2 of Fr\'ed\'eric Goualard, and most programs
written for GAOL~4 compile and run with it once five things are done: the
namespace is opened, the strings they gave to \code{interval(const char*)}
are read by \code{textToInterval()}, the relations \code{==} and \code{!=} on
intervals are rewritten with \code{set\_eq()}, \code{set\_disjoint()} or
\code{!set\_disjoint()}, the handlers of the exceptions of GAOL~4 are
retargeted on those of the standard library, and the conversions of a
\code{double} or an integer to an interval, explicit now, are written out.
The other changes concern what a program
sees: tighter bounds, the empty set where GAOL~4 built an interval, the
comparisons of IEEE 1788-2015 \cite{IEEE1788:2015}, and texts read and written
otherwise. Each change was made for a reason, and most of them were measured.

\medskip
In this chapter, we gather what a program may have to change and what it
sees differently. The first section lists what concerns the compilation and
the link, the second the results that change. The third section is the
register of the differences, theme by theme: for each change, what GAOL~4
did, what GAOL v5 does, why, and what a program sees, with the measurements
that justified it and the chapter of this manual that describes the new
behaviour. The last section discusses the choices made.

\medskip
A word on the versions. ``GAOL~4'' is the 4.2.2, which GAOL v5 continues;
the benchmarks of Chapter~\ref{chap:computing} use the 4.2.3, the last
version of Fr\'ed\'eric Goualard, which came after it. The releases 4.3.0 to
4.3.2 of the repository of GAOL v5 came next and are not GAOL~4: they already
had \code{rnd\_keep()}, \code{get\_rounding()} and \code{MEMFREE()}
(\code{GAOL\_MEMFREE()} now), which neither the 4.2.2 nor the 4.2.3 has. The
meson build that this chapter counts with GAOL~4 is the one added to the
repository of the 4.2.3 in November 2025, after its last release: the 4.2.2 had
\cmd{configure} only.
GAOL~4 also had classes of exceptions of its own,
\code{input\_format\_error}, \code{invalid\_action\_error} and
\code{unavailable\_feature\_error}, derived from \code{gaol\_exception}. GAOL
v5 throws the \code{std::invalid\_argument} and the \code{std::domain\_error}
of the standard library where GAOL~4 threw the first two, and has none of
these classes (\pxref{sec:exceptions}). Where this chapter says what GAOL~4
threw, it gives the class of GAOL~4.

\section{Compiling and linking}\label{sec:migration-compile}
%%-----------------------------
\cindex{GAOL~4!compiling a program of}%

The changes below are the ones a compiler or a linker reports, and those of
the initialization and the cleanup, which a program makes in its text. Most of
them are described in the section the cross-reference names.

\begin{itemize}
\item[$\bullet$] \textbf{The namespace.} \file{gaol/gaol.h} no longer opens
  the namespace \code{gaol}: a program adds \code{using namespace gaol;}, or
  names \code{gaol::interval} (\pxref{sec:namespaces}). The type and the
  functions are in \code{gaol\_core}, and the symbols of the library name it:
  a program compiled against the headers of GAOL~4 has to be compiled again.
\item[$\bullet$] \textbf{The names GAOL keeps for itself.} \code{gaol} holds
  the names GAOL~4 had there that GAOL v5 still has, except those GAOL uses
  for itself, which a program names in \code{gaol\_core}
  (\pxref{sec:namespaces}): \code{gaol\_signbit()}, \code{uintdouble},
  \code{ullidouble}, \code{Interval\_struct}, \code{the\_null\_expr},
  \code{prec\_t}, \code{modulo\_k\_pi()} and the \code{round\_*\_sse()}, all of
  them internals that the manual of GAOL~4 did not document, and, of the 4.3
  releases of GAOL v5,
  \code{rounding\_state}, \code{get\_rounding()}, \code{set\_rounding()},
  \code{round\_upward\_if\_needed()} and \code{gaol\_set\_rounding\_x86()}.
\item[$\bullet$] \textbf{Visual C++.} The code using GAOL need no longer define
  \code{\_\_GAOL\_PUBLIC\_\_} empty: GAOL is a static library there, and its
  declarations carry no \code{\_\_declspec(dllimport)}. A definition left in
  a project changes nothing. The macro is \code{GAOL\_PUBLIC} now.
\item[$\bullet$] \textbf{The allocator.} The allocator of containers
  \code{aligned\_allocator} is gone, with its header,
  \file{gaol/gaol\_allocator.h}.
\item[$\bullet$] \textbf{One library to link.} GAOL is the only library to
  link: the mathematical library GAOL~4 was built with is no longer used,
  CORE-MATH being compiled into GAOL, and the options of the builds that chose
  it are gone.
\item[$\bullet$] \textbf{The flags of interval arithmetic.} The code
  including the headers of GAOL is compiled with the flags of interval
  arithmetic, which \code{gaol::gaol} and \file{gaol.pc} give
  (\pxref{sec:compiling}); \file{gaol/gaol\_config.h} refuses the options and
  the compilers that give wrong bounds (\pxref{sec:refused}).
\item[$\bullet$] \textbf{The intervals of floats.} \code{gaol::intervalf} and
  \code{gaol::interval2f} are compiled by none of the builds, which have no
  option for them: only a developer of GAOL compiles them, defining
  \code{GAOL\_FLOAT\_INTERVALS} in \file{gaol/gaol\_config.h}.
\item[$\bullet$] \textbf{The bounds at doubles.} The functions of doubles
  rounded downward and upward, \code{exp\_dn()}, \code{sin\_up()},
  \code{ipow\_up()}, \code{ipow\_dn()}, \code{pow\_up()}, \code{pow\_dn()}
  and the others of \file{gaol/gaol\_double\_op.h}, are no longer part of the
  interface: the header is no longer installed. \code{uipow\_upup()} and
  \code{uipow\_dnup()} are no longer declared.
\item[$\bullet$] \textbf{The static constants.} The static intervals
  \code{interval::cst\_one}, \code{cst\_pi} and the other \code{cst\_}
  members are gone: \code{interval::one()}, \code{interval::pi()} and the
  other functions give these constants, right even in the initialization of a
  static object (\pxref{chap:init}).
\item[$\bullet$] \textbf{The initialization.} \code{gaol::init()} need not be
  called again after the program changed the rounding direction: each
  operation sets it upward itself (\pxref{sec:rounding-direction}). It need
  not be called at all, GAOL initializing itself with every compiler, Visual
  C++ included (\pxref{chap:init}).
\item[$\bullet$] \textbf{The cleanup.} The program calls
  \code{gaol::cleanup()} right after its last use of GAOL: unless the
  rounding direction is preserved, it sets back the rounding direction that
  the initialization found, which GAOL~4 left upward when a static object of the
  program had computed before its initialization (\pxref{chap:init}). It
  no longer frees what the initialization allocated, which the automatic
  cleanup frees when the program ends.
\item[$\bullet$] \textbf{The relations.} The operators \code{==} and
  \code{!=} on intervals are gone, with the possibly relations
  (\code{possibly\_eq()}, \code{possibly\_le()}\ldots),
  \code{certainly\_eq()} and \code{certainly\_neq()}, and so is the option
  \loption{enable-relations} of \cmd{configure} and meson: a program compares
  intervals with \code{set\_eq()}, \code{set\_disjoint()} and
  \code{!set\_disjoint()}, and \code{<}, \code{<=}, \code{>} and \code{>=}
  are the certainly relations (Chapter~\ref{chap:relations}). The new
  relations do not give the results of the old operators: \code{==} was
  \code{certainly\_eq()} in every build by default, true for two empty sets and
  for a point interval equal to the upper bound of the other interval
  (\code{[4] == [3, 4]} was true, \code{[3, 4] == [4]} false), and \code{!=} was
  its negation, so that
  \code{==} of $[3, 4]$ and $[3, 4]$ was false, where \code{set\_eq()} is
  true, and \code{!=} was true, where \code{set\_disjoint()} is false. A
  program chooses the relation it means.
\item[$\bullet$] \textbf{The constructors from strings.} The constructor from
  one string, \code{interval(const char*)}, is gone, and with it the
  conversion of a string wherever an interval is expected:
  \code{textToInterval("0.1")} reads a string in its place, and throws
  \code{std::invalid\_argument} where the constructor threw
  \code{input\_format\_error}; \code{x = "[1, 2]"} becomes \code{x = textToInterval("[1, 2]")}, and
  \code{x + "0.1"} becomes \code{x + textToInterval("0.1")}
  (\pxref{sec:constructors}). The constructor from two strings is gone too:
  \code{interval("0.1", "0.3")} becomes \code{textToInterval("0.1", "0.3")},
  which gives the empty set for bounds in the wrong order. \code{interval(0)}
  and \code{interval(0, 0)}, which these constructors made ambiguous, compile.
\item[$\bullet$] \textbf{The constructor from a double.}
  \code{interval(double)} is explicit, and so is the constructor from an
  integer: \code{interval x = 0.5;}, \code{interval x = 5;},
  \code{return 0.5;} in a function returning an interval, \code{min(x, 0.5)}
  and \code{interval r = sqrt(2.0);} do not compile, where GAOL~4 converted the
  double without a word. \code{x = 0.5} and \code{x < 0.0}
  still compile (\pxref{sec:constructors}).
\item[$\bullet$] \textbf{The exceptions.} The classes \code{gaol\_exception},
  \code{input\_format\_error}, \code{invalid\_action\_error} and
  \code{unavailable\_feature\_error} are gone, with
  \file{gaol/gaol\_exceptions.h}: a handler names
  \code{std::invalid\_argument} or \code{std::domain\_error}
  (\pxref{sec:exceptions}).
\item[$\bullet$] \textbf{The headers.} \file{gaol/\allowbreak gaol\_profile.h},
  \file{gaol/\allowbreak gaol\_parameters.h}, \file{gaol/\allowbreak gaol\_limits.h},
  \file{gaol/\allowbreak gaol\_assert.h}, \file{gaol/\allowbreak gaol\_exact.h} and
  \file{gaol/\allowbreak gaol\_flags.h} are gone, and so are
  \file{gaol/\allowbreak gaol\_config\_msvc.h},
  \file{gaol/\allowbreak gaol\_config\_mingw.h},
  \file{gaol/\allowbreak gaol\_version\_msvc.h},
  \file{gaol/\allowbreak gaol\_version\_mingw.h} and
  \file{gaol/\allowbreak gaol\_flt\_output.h} (the source
  \file{gaol/\allowbreak dllmain.cpp} too); \file{gaol/\allowbreak gaol\_init\_cleanup.h} is no longer
  installed, and a program that builds expressions includes
  \file{gaol/parser/\allowbreak gaol\_expression.h}, with
  \file{gaol/parser/\allowbreak gaol\_expr\_eval.h} to evaluate them, where it
  included \file{gaol/\allowbreak gaol\_expression.h} (\pxref{sec:changes-interface}).
\item[$\bullet$] \textbf{The expressions.} \code{eval\_stack} and
  \code{expr\_eval} are no longer copyable, the stack of an evaluation
  growing in the heap (\pxref{sec:changes-io}).
\item[$\bullet$] \textbf{The debugging macro.} \code{GAOL\_DEBUG(level,
  command)} of the debug build is \code{GAOL\_DEBUG\_VERBOSE}: the old name is
  not kept (\pxref{sec:debugging-facilities}).
\item[$\bullet$] \textbf{The options of the builds.} The option
  \code{with-test} of meson is \code{with-tests}. \loption{enable-asm} of
  \cmd{configure} and \code{enable-asm} of meson are gone, GAOL v5 no longer
  calling the assembly they chose: \cmd{configure} warns of an option it does
  not know, and \cmd{meson setup -Denable-asm=true} stops with an error
  (\pxref{sec:changes-builds}).
\end{itemize}

\section{Results that change}\label{sec:migration-results}
%%---------------------------
\cindex{GAOL~4!results that change}%

The changes below are the ones a program sees at run time without any
change of its text. The third section gives their reasons.

\begin{itemize}
\item[$\bullet$] \textbf{Static objects.} The intervals computed in the
  initialization of the static objects of a program linked with the static
  library GAOL are right: \code{interval::pi()} was $[-0, 0]$ there, and
  $[1, 2]/[-1, 1]$ was $[-0, 0]$ rather than $[-\infty, +\infty]$, the
  constants of GAOL being computed after these objects, except with MinGW-w64
  (\pxref{chap:init}).
\item[$\bullet$] \textbf{The bounds are tighter.} The elementary functions
  give the tightest bounds, rather than one double beyond, and are exact
  where their value is a double; the integer powers are
  computed from exact products, and the $n$th roots are proved with powers
  rounded upward and downward; \code{sin}, \code{cos} and \code{tan} are
  within one double of the tightest bounds at every magnitude
  (Chapter~\ref{chap:accuracy}).
\item[$\bullet$] \textbf{The empty set.} The constructors give it for
  \itv{+\infty}{+\infty}, \itv{-\infty}{-\infty}, bounds in the wrong order
  and NaN bounds, and assigning an infinity to an interval empties it
  (\pxref{sec:constructors}); the intersection of disjoint intervals is it,
  with NaN bounds; so are \code{log} of an interval holding no positive
  number, \code{invabs\_rel(J, I)} for a negative \var{J} and \code{div\_rel} with
  no solution in \var{I}. The operators of an interval and an infinite
  \code{double} (\code{x * INFINITY}, \code{x + INFINITY}) give it too.
  \code{width()} of the empty set is NaN rather than $-1$.
\item[$\bullet$] \textbf{Powers.} \code{pow(x, p)} for a double \var{p} no
  longer truncates \var{p}; \code{pow(x, y)} keeps the negative part of
  \var{x} for an integer exponent only, and takes the \ieee{pow} of IEEE
  1788-2015 otherwise (\pxref{sec:powers}). \code{nth\_root} of an odd order
  is defined on the whole real line, and takes a negative order.
\item[$\bullet$] \textbf{The comparisons} follow IEEE 1788-2015 for the empty
  set and the infinite bounds: \code{certainly\_leq()} and its companions are
  true when either interval is empty, and \code{set\_strictly\_contains()}
  takes an infinite bound as beyond the same infinite bound (Chapter~\ref{chap:relations}).
\item[$\bullet$] \textbf{Input and output.} Numbers are read exactly, and the
  literals of IEEE 1788-2015 are read, in any case of letters;
  \code{"[inf]"} is the empty set. Every string is an expression, where an
  interval may stand wherever a number may (\code{"1+[1,2]"}), an uncertain
  number being no bound (\code{"[5?1]"} is refused) and needing a blank
  after a sign next to it (\code{"10-3.56?1"} is refused,
  \code{"10 - 3.56?1"} is read), and an error anywhere
  stops the reading. The reader knows every function of GAOL, and
  \code{gaol\_ieee1788::textToInterval} reads the names of IEEE 1788-2015
  (\code{pown}, \code{rootn}\ldots). The format \code{hexa} writes interval
  literals, which are read back bit for bit; the decimal formats round the
  bounds outward whatever the C library, and the format \code{agreeing} keeps
  the digits of each bound (Chapter~\ref{chap:io}).
\item[$\bullet$] \textbf{New operations.} \code{atan2()} is implemented, and
  the functions of IEEE 1788-2015 GAOL did not provide are added
  (\code{exp2}, \code{hypot}, \code{sign}, \code{trunc}\ldots,
  Chapter~\ref{chap:arithmetic}).
\end{itemize}

\section{The differences, by theme}\label{sec:migration-themes}
%%--------------------------------------------
\cindex{GAOL~4!differences with}%

This section condenses the register of the differences between GAOL~4 and
GAOL v5, theme by theme. The themes are those of the register, regrouped: the
elementary functions first, then the rounding direction, the floating-point
exceptions, the constants, the kernels, the intervals and their functions, the
powers and the roots, the other functions, the relations, the reading and the
writing, the errors, the interface and the builds. Some of the defects told
below were those of the earlier versions of GAOL v5 only, and the text says
so. Each measurement is given with the machine and the compiler it was made
on wherever the register gives them. The timings are those of the step that
made each change, which may be an intermediate one; those of the release
against GAOL 4.2.3 are in Chapter~\ref{chap:computing}.

\subsection{The elementary functions and CORE-MATH}
\label{sec:changes-core-math}
\cindex{CORE-MATH!differences with GAOL~4}%

Every elementary function is bounded with CORE-MATH \cite{CoreMath:2022}, on
every architecture and with every compiler, and it is the only library GAOL
takes its elementary functions from.\newinvfive{} The other mathematical
libraries GAOL~4 could be built with are gone, with their sources, their
headers and the scripts that installed them, and so are the options of the
builds that chose one: there is nothing to choose, to find, to build apart or
to link along with GAOL, whose library holds the functions
(\pxref{sec:math-libraries}). No build names \option{-lm}: the C library
still gives GAOL \code{sqrt()}, \code{floor()}, the functions of
\file{<fenv.h>} and, where the processor has no fused multiply-add
instruction, \code{fma()}, but \cmd{g++} and \cmd{clang++} link the math
library themselves, so that the autotools build (whose \code{LDADD} was
\option{-lm} in GAOL~4), meson and the \file{gaol.pc} the builds install do
not name it, as the CMake build never did.

\medskip
The bounds are the tightest ones rather than one double beyond. CORE-MATH is
correctly rounded in the rounding direction in effect, so GAOL takes the
value at a bound, computed in the upward rounding it keeps, as the upper
bound, and the double below it as the lower one. GAOL~4 took values correctly
rounded to nearest only, so each bound had to be moved one double outward,
and the direction had to be set to nearest and back around every call. On an
Intel i7-1185G7 with Clang 18, the bounds of \code{log} took 30 ns rather
than 62. The accurate phases of \code{log},
\code{sin}, \code{cos}, \code{tan}, \code{atan2} and \code{pow} compute with
a 128-bit integer, which Visual C++ has on no architecture and GCC has on no
32-bit target: \file{gaol/math/gaol\_u128.h} gives them the type of the
compiler where it has one and a structure of two 64-bit halves where it has
none. The two paths were compared with the pristine upstream sources over
about 40 million arguments in the four rounding directions, and give the same
bits; a job of the continuous integration builds the whole of GAOL with the
halves forced (\code{-DGAOL\_U128\_EMULATION=ON}, in \file{linux.yml};
\pxref{sec:ci}).

\medskip
MinGW-w64 is no longer refused for its version alone, both former reasons
being gone: the hyperbolic functions of its math library, which GAOL no
longer uses, and the cost of its \code{fesetround()}, which GAOL no longer
calls for its elementary functions.\newinvfive{} Still refused are the MinGW-w64 whose
\code{fma()} or \code{round()} is wrong, those of its own math library,
computed in doubles: that \code{fma()} is not correctly rounded (10.7\,\% of
the error-free products \code{fma(a, b, -a*b)} were wrong with mingw-w64 11
on x86-64), and CORE-MATH computes with it wherever the compiler has no
\option{-mfma}; that \code{round()} depends on the rounding direction
(\pxref{sec:refused}).

\medskip
The functions of IEEE 1788-2015 that GAOL~4 did not provide are
added.\newinvfive{} \code{exp2}, \code{exp10}, \code{log2} and \code{log10},
required among the forward elementary functions (Table~9.1 of the standard),
are bounded with CORE-MATH's, the tightest bounds, and the value that is a
double is kept as a bound rather than moved: $2^3 = 8$, $10^{22}$,
$\log_2(1/4) = -2$ and $\log_{10}(100) = 2$ are exact. Sixteen functions of
Table~10.5, which the standard recommends, follow: \code{expm1},
\code{exp2m1}, \code{exp10m1}, \code{log1p}, \code{log2p1}, \code{log10p1},
\code{hypot}, \code{rsqrt}, \code{sinpi}, \code{cospi}, \code{tanpi},
\code{asinpi}, \code{acospi}, \code{atanpi}, \code{atan2pi} and
\code{compoundm1}. \code{sinpi}, \code{cospi} and \code{tanpi} find their
extrema and poles exactly, at the multiples of $1/2$, which are doubles:
\code{sinpi([1e17])} is 0, where \code{sin(pi*x)} gave \itv{-1}{1}, the
width of $\pi$ times $10^{17}$ being more than a period.
$\mathrm{compoundm1}(x, y) = (1 + x)^y - 1$ is the one CORE-MATH has no
binary64 version of, and $\exp(y\,\mathrm{log1p}(x)) - 1$ computed with
doubles is up
to 1\,100 doubles from the value: it is computed with integers only, in
64-bit fixed point where that proves the bounds (95\,\% of the pairs for
$|y\,\mathrm{log1p}(x)|$ below 0.35, in 0.1~$\mu$s a corner) and with a 128-bit
significand in balls otherwise (0.6 to 0.9~$\mu$s). \ieee{fma},
\ieee{cancelMinus} and \ieee{cancelPlus} (\code{cancel\_minus} and
\code{cancel\_plus}), which the standard requires the tightest, and
\ieee{less}, \ieee{strictLess}, \ieee{isEntire} and \ieee{isCommonInterval}
are provided too. The error function \code{erf} and
its complement \code{erfc}, which the standard does not name, are bounded
with CORE-MATH's as well, exact at 0 and at the infinities. The integer
functions \code{sign}, \code{trunc}, \code{round\_ties\_to\_even} and
\code{round\_ties\_to\_away} are exact, each being non-decreasing; they take
nothing from CORE-MATH, and \code{round\_ties\_to\_even} reads the bits so as
to ignore the rounding direction, which \code{std::nearbyint} and
\code{std::rint} read (Chapters~\ref{chap:arithmetic} and
\ref{chap:functions}). The namespace \code{gaol\_ieee1788} gives the
operations of the standard GAOL provides, under the names and in the argument
order of the standard (Chapter~\ref{chap:ieee1788}).

\medskip
CORE-MATH reads the rounding direction as GAOL sets it.\newinvfive{} GAOL~4
had no CORE-MATH, and the defects below were found and fixed in the earlier
versions of GAOL v5 that took it. Some
sources of CORE-MATH round a result themselves, in the direction they read,
and two of the ways they read it gave another direction than that of the
doubles GAOL computes, the one of MXCSR. \code{pow} rounds a subnormal
result in the direction \code{fegetround()} gives, which the C library of
Linux reads from the x87 unit: with the x87 unit to nearest and MXCSR upward,
the upper bound of \code{pow} was below the exact value for about half of the
arguments with a subnormal result. GAOL gives the sources a
\code{fegetround()} that reads MXCSR on x86. On 32-bit x86 that spares
\code{cbrt} and \code{asinpi} a call of the C library at each bound, about
80 ns a call with Visual C++ for x86: on the runners of the continuous
integration, \code{nth\_root(x, 3)} takes 179 ns rather than 355 and
\code{asinpi(x)} 794 ns rather than 1005 with Visual C++ for x86, and 120 ns
rather than 324 and 255 ns rather than 438 with clang-cl for x86.
\code{cbrt}, \code{rsqrt} and \code{asinpi} read the rounding field of MXCSR
and turned it into the \code{FE\_*} value of the C library by a shift chosen
from the macros of the platform; with clang-cl and with Cygwin on x86-64,
whose newlib gives \code{FE\_*} the values 0 to 3, the shift was wrong for
the three, and the directed roundings were taken for toward zero
(\code{asinpi} was wrong for about half of its arguments with clang-cl). The
field is 0 to nearest, 1 downward, 2 upward and 3 toward zero on every x86
processor, and the three functions now take it so, whatever the values of
the C library.

\medskip
CORE-MATH's \code{pow} is correctly rounded in every direction. A test of
\code{pow} where $x^y$ is a double found three defects in the earlier
versions of GAOL v5. With Visual C++,
\code{NAN} and \code{INFINITY}, which the UCRT writes for C as products of
\code{1e+300}, were computed when the program runs, in the rounding direction
in effect, so that \code{pow} rounded downward was $-0$ from $2^{-1075}$ to
about $2^{-947}$; they are read from their bits now. On 32-bit x86 the square
root of the C libraries of Windows rounds to nearest in every rounding
direction: CORE-MATH takes its square root from the SSE2 instructions there.
On 32-bit ARM, the conversion of a 54-bit integer rounded to nearest, so that
the upper bound fell below $x^y$ at a midpoint. Moreover, for $y \log x$
from 709.78263 to 709.78267, a power between $1.7975 \cdot 10^{308}$ and
$1.7976 \cdot 10^{308}$, the first phase of \code{pow} computed
$\infty - \infty$ before its second phase gave the right result, and a
program that had enabled the invalid-operation exception died on SIGFPE: the
vendored source sends those arguments to the second phase. The source of
\code{cospi} shifted a signed integer out of its range, undefined behaviour
that the tests run with UBSan reported; it shifts it unsigned now.

\medskip
The cube root \code{nth\_root(x, 3)} is CORE-MATH's \code{cbrt}, the
tightest bounds, exact where the cube root is a double, which cubing the
value tells: it takes 41 ns, where the search by bisection the other roots
use takes 89 (the medians of three runs, Intel i7-1185G7, Clang 18.1; 24 ns
against 77 on an Apple M2, Apple clang 21).\newinvfive{} \code{nth\_root(x, q)} takes
a negative $q$, which the standard recommends (\ieee{rootn} over
$\mathbb{Z}\setminus\{0\}$): the reader of strings converted the exponent to
an unsigned int, and \code{nth\_root(16, -2)} read from a string was the
4294967294-th root of 16 rather than \itv{0.25}{0.25}. The second argument of
\code{nth\_root}, \code{pown} and \code{rootn} read in a string may be any
integer a double holds: \code{nth\_root(8,3000000000)} threw
\code{invalid\_action\_error}, and \code{rootn(8,3000000000)} read with the
names of the standard was the empty set in the earlier versions of GAOL v5,
where the same call from C++ is computed. The reader of strings also takes the new names \code{exp2},
\code{log2}, \code{cbrt}, \code{sign} and \code{trunc}, which gave a syntax
error before (\pxref{sec:input-format}).

\subsection{The rounding direction and the initialization}
\label{sec:changes-rounding}
\cindex{rounding direction!differences with GAOL~4}%

The rounding direction is set upward by each operation that computes a bound
when it is not, in every build; the operations of the AVX-512 path set none,
their instructions carrying their rounding.\newinvfive{} An operation that
rounds nothing (\code{-x}, \code{abs()}, \code{min()}, \code{floor()}, the
predicates) or returns before its check (\code{y * empty},
\code{exp(empty)}) leaves the direction as it found it
(\pxref{sec:rounding-direction}). GAOL~4 has no CMake build; the one of the 4.3
releases of GAOL v5, which derives from that of IBEX, set it once, in
\code{gaol::init()}, and computed in the direction the calling code left:
rounding to nearest, downward or toward zero, 30 operations gave other
bounds, and sums and products did not enclose their exact values (3199 of the
13825 checks of the test of the rounding direction failed with GCC 9.4 on
x86-64). Built with \cmd{configure} or meson, whose option
\loption{enable-preserve-rounding} is on by default in GAOL~4 although its
manual says the contrary, GAOL~4 restored the rounding direction after each
operation, but only the x87 control word on x86-64,
where GAOL computes with SSE (43 operations left the SSE instructions
rounding to nearest or upward, whatever direction they found), and without
the rounding bits of the FPCR register on arm64.

\medskip
With the option \code{GAOL\_PRESERVE\_ROUNDING} of the three builds
(\loption{enable-\allowbreak preserve-\allowbreak rounding} of \cmd{configure} and
\option{-Denable-preserve-rounding} of meson, which GAOL~4 had),
GAOL saves and restores the
direction by reading and writing the control registers itself, a register
only if it differs from what is set back: the x87 control word and MXCSR of
an x86 processor, FPCR with the NEON intervals, FPSCR with the FPU intervals
of 32-bit ARM processors and GCC; the other processors keep
\code{fegetround()} and \code{fesetround()}. Each value computed before the
rounding direction changes is written to a \code{volatile} variable: GCC does
not implement \code{\#pragma STDC FENV\_ACCESS}, and moved the computation of
the midpoint after \code{fesetround()} (GCC bug 34678).\newinvfive{} Whether the
direction is upward is shown by an addition, $1 + 2^{-60}$ with $2^{-60}$
read from \code{volatile} memory, wherever doubles are computed in double
precision; on x86-64 processors and, with GCC and Clang, on ARM ones it is
$1 + (2^{-1060} + 0)$, which shows the modes that flush the subnormal numbers
to zero as well. The NEON intervals of 64-bit ARM processors read the
register FPCR instead, and a 32-bit x86 processor the control registers of
the two units that round. Reading the direction with the SSE register and
\code{fegetround()} rather than by an addition made \code{x + y} take 800 ns
rather than 2.8 ns under Rosetta 2, 82 to 98 ns rather than 2.6 to 4 ns with
32-bit Visual C++, and 48 ns rather than 6.4 ns with 32-bit MinGW-w64 15.2;
and Clang 18 read MXCSR once for a whole loop that changed the rounding
direction.

\medskip
The direction is checked once per operation, but for two cases told
below.\newinvfive{} The bounds of the
elementary functions at doubles are computed by functions that take the
direction to be upward already, as it is after the check of the function of
intervals calling them: in the earlier versions of GAOL v5, \code{exp()} of an
interval checked the direction three times, \code{sin()} and \code{cos()} up to four, the other elementary
functions and \code{pow(x, y)} two or three. \code{exp()} takes 5.7\,\% less
time, \code{sin()} 3.6\,\%, \code{pow(x, y)} 3.4\,\% (Clang 18, Intel
i7-1185G7). The operations calling others check once too: in those versions,
\code{tan()} checked six times, calling \code{width()}, \code{+} and \code{/} on
intervals after its own check, and \code{acos\_rel()}, \code{asin\_rel()}
and \code{atan\_rel()} six times too; they compute on doubles after their one
check, and give the same bounds, bit for bit, but for the quotients of
\code{modulo\_k\_pi()}, now computed on doubles. \code{tan()} takes 22\,\%
less time, \code{acos\_rel()} 24\,\%, \code{sqrt\_rel()} 29\,\%,
\code{pow(x, [3])} 7\,\%, \code{nth\_root(x, -2)} 13\,\%, and
\code{modulo\_k\_pi()} 10 ns rather than 22; with
\code{GAOL\_PRESERVE\_ROUNDING}, where each check saves and restores the
direction, \code{tan()} 74\,\%, \code{acos\_rel()} 83\,\%,
\code{pow(x, [-3])} 20\,\% and \code{nth\_root(x, -3)} 15\,\% (Clang 18,
Intel i7-1185G7). Still checking more than once are \code{nth\_root(x, q)}
for an order beyond 64 bits, which only an \code{\_\_int128} gives, and
\code{nth\_root(x, -2)} where the square root has a vector path of its own (the
AVX-512 path, taken with \code{GAOL\_PRESERVE\_ROUNDING}, and the NEON intervals
without it), which takes \code{sqrt()} and then \code{inverse()}, each
checking, faster than one check before the scalar root. With that option,
\code{GAOL\_RND\_KEEP()} of an interval
keeps its two bounds through two volatile doubles with the SSE2 and FPU
intervals and in the register itself with the NEON ones, where the 4.3
releases copied the interval a byte at a time into volatile memory: \code{pow(x, [-3])} and
\code{nth\_root\_rel()} take 26 to 27\,\% less time, \code{asinh\_rel()} and
\code{atanh\_rel()} 21\,\% (Clang 18, Intel i7-1185G7). \code{tan()} and
\code{pow(x, y)} keep their bounds before setting the direction back, which
they did not, so that GCC, which does not model the rounding direction,
could compute them after; and the functions of doubles \code{exp\_dn()},
\code{sin\_up()} and the thirty others saved the direction they found but
left it upward, a defect the compilers only showed by warning that a variable
was set but not used. They no longer save the direction, and their header,
\file{gaol/math/gaol\_double\_op.h}, is no longer installed.

\medskip
The initialization and the cleanup change.\newinvfive{} GAOL initializes
itself with every compiler, before the static objects of the program: each
file including \file{gaol/gaol\_common.h} holds a static
\code{gaol\_core::gaol\_initializer}, whose constructor initializes GAOL the
first time, as \code{std::ios\_base::Init} does the standard streams. GAOL~4
declared \code{gaol\_init()} with \code{\_\_attribute\_\_((constructor))},
for GCC and Clang only: Visual C++ had no automatic initialization, and the
constructor of the static library ran after those of the program (except with
MinGW-w64, which runs the constructors from the last linked to the first), so
that the static objects of the program computing intervals set the rounding
direction upward before the first \code{gaol::init()} read it, and
\code{gaol::cleanup()} left it upward. \code{gaol\_init()} is gone.
\code{gaol::cleanup()} sets back the rounding direction that the first
initialization found, unless it is preserved: GAOL~4 left it upward after its
use when a static object of the program had computed before the
initialization, even after its automatic cleanup, and a program computing to
nearest after it had to set it back itself. \code{gaol::cleanup()} does so at its
first call only, and an operation of GAOL after it sets the direction upward
again; \code{gaol::restore\_rounding()} sets it back as many times as the
program needs. Only the rounding direction is set back, not the rest of the
environment: the exception flags raised since are the program's.
\code{gaol::cleanup()} no longer frees what the initialization allocated: it
deleted the node of the empty expression while expressions still referred to
it, a local of \code{main()} destroyed after the call or a static object,
whose destructors then wrote into freed memory. The automatic cleanup frees
it when the program ends or the library is unloaded (\pxref{chap:init}).
Following the patch IBEX applies to GAOL, the FPU is initialized with
\code{fesetenv(FE\_DFL\_ENV)} and \code{round\_upward()}, rather than by
writing an x87 control word into \code{fenv\_t}, which on ARM64 left the
rounding to nearest and crashed an ARM64 Mac; and \code{interval::midpoint()}
and \code{operator<<} set the rounding direction back upward, which they did
not when the rounding direction is not preserved.

\medskip
The rounding direction is set on x86 processors by writing the control
registers of the x87 and SSE units (\code{fnstcw}/\code{fldcw},
\code{stmxcsr}/\code{ldmxcsr}) rather than through \code{fesetround()},
which cost 130 ns per call with mingw-w64 13, 50 ns with the C runtime of
Visual C++ for x64 and 250 ns for x86, 8.5 ns with glibc.\newinvfive{} GAOL~4
changed the direction four times for each elementary function of an interval,
where the math library itself took about 10 ns; GAOL v5 no longer changes it
there, CORE-MATH computing upward. On the runners of the continuous
integration, \code{exp()}, \code{log()}, \code{sin()} and \code{cos()} of an
interval took 535 to 630 ns with MinGW-w64 and MSYS2, and take 48 to 209 ns;
250 to 340 ns with Visual C++ for x64, and 45 to 160 ns; 1050 to 1750 ns with
Visual C++ for x86, and 150 to 650 ns. A register is written only if its
direction differs from the one set, and no other bit changes. With mingw-w64
13 and later there is a second reason: its \code{fesetround()} clears the six
exception flags of MXCSR when it changes the register, and its
\code{fegetround()} joins by OR the rounding fields of the two units. On a
32-bit x86 processor GAOL also reads the two registers in line, in one asm
statement, rather than through \code{fegetround()}, about 80 ns with Visual
C++ for x86. With \code{GAOL\_PRESERVE\_ROUNDING}, \code{x.midpoint()} takes
17.5 ns rather than 27.9 (GCC 15.2, Intel i7-1185G7), and 52 ns rather than
131 with Visual C++ for x64 on the runners of the continuous integration;
the bounds are the same. The option also writes the registers only if they
have to change: an operation of the SSE2 intervals read MXCSR and wrote it
twice whatever it held. In a program that computes upward, \code{x + y}
takes 1.7 ns rather than 8.2 and \code{x * y} 2.5 ns rather than 8.3; in a
program that computes to nearest, 0.2 ns more (GCC 15.2, Intel i7-1185G7, the
best of 15 runs of a loop). Visual C++ for x64 writes MXCSR alone with the
option, as GCC, Clang and clang-cl do with the fused multiply-add
instructions: in the earlier versions of GAOL v5, the operations that start
with \code{GAOL\_RND\_ENTER()} called the C runtime, because the macros asked for \code{\_\_FMA\_\_}, which
\cmd{cl} does not define; they take \code{\_\_AVX2\_\_} too now, and
\code{pow(x, 3)} takes 29 ns rather than 92, \code{sin(x)} 71 ns rather than
135 on the runners of the continuous integration (AMD EPYC, Visual Studio
2026).

\medskip
On ARM processors GAOL writes the control register itself too.\newinvfive{} With the NEON
intervals of 64-bit ARM processors the direction is the field RMode of FPCR,
read with \code{mrs} and written with \code{msr} (\pxref{sec:neon}). On
32-bit ARM processors with GCC and a floating-point unit of doubles it is the
field RMode of FPSCR, read with \code{vmrs} and written with \code{vmsr}:
\code{x + y} takes 22.5 ns rather than 49 with \code{GAOL\_PRESERVE\_ROUNDING}
(the Debian armhf jobs of the continuous integration, which run on the arm64
runners of GitHub). FPSCR holds the flags of the exceptions too, so every
write but the one at the start of an operation reads the register again just
before, and no flag raised since is lost. Without the option nothing changes
there: the check of each operation stays the addition, \code{vmrs} costing
more than its two additions (\code{x + y} 6.1 ns rather than 4.1). Elsewhere
(POWER, s390x, RISC-V, MIPS, the FPU intervals of 64-bit ARM, and 32-bit ARM
without GCC or without a floating-point unit of doubles), \code{fesetround()}
still.

\subsection{The floating-point exceptions and the flush-to-zero modes}
\label{sec:changes-exceptions}
\cindex{floating-point exceptions!differences with GAOL~4}%

Operations raise no invalid-operation exception for the empty set, nor for a
NaN the program gives, nor for a zero bound times an infinite one in an SSE2
register.\newinvfive{} The bounds of the empty set are NaN, and
\code{is\_empty()} was \code{!(left() <= right())}, as \code{set\_contains()},
\code{set\_disjoint()} and the other relations compared them with \code{<},
\code{<=}, \code{>} and \code{>=}: comparisons that signal. With the
exception enabled, \code{is\_empty()} of an empty interval and the operations
that call it stopped the program with SIGFPE, and with it masked they raised
the flag. The SSE2 product computed $0 \cdot \infty$ and cleared the NaN
afterwards, and the SSE2 \code{div\_rel} divided an infinite bound by an
infinite one in the half of a register it throws away. The bounds of an empty
operand are now compared with the quiet comparisons of \file{<cmath>}
(\code{std::islessequal()}\ldots), one instruction with GCC and Clang on
x86-64, as is the first comparison of the constructor, which gives the empty
set for NaN bounds: every operation, relation, function and output takes the
empty set without the exception, and so does a NaN the program gives
(\code{interval(NAN)}, \code{x += NAN}, \code{x <= NAN}, \code{pow(x, NAN)}).
The SSE2 divisions, \code{\%}, \code{inverse()} and \code{div\_rel()} raise
no division-by-zero exception for a divisor with a zero bound
(\code{1/[0, 1]}, \code{inverse([0, 2])}), which they raised in GAOL~4, and
\code{midpoint()} of \code{[DBL\_MAX]} raises no overflow exception, a bound
of $2^{1023}$ or more in magnitude telling the case of the sum of the halves
before the bounds are added. What stays: the inexact exception is raised by
almost every operation, and \code{log([0, 1])} and \code{[1e308] * 10}, whose
bounds are infinite, still raise the division-by-zero and the overflow ones.
A program that enables an exception has to disable it while GAOL computes
(\pxref{sec:fp-exceptions}).

\medskip
A result with no infinite bound raises no overflow nor division-by-zero
exception.\newinvfive{} Some functions computed beyond the largest double on
the way, or divided by zero, and raised \code{FE\_OVERFLOW} or
\code{FE\_DIVBYZERO} for a finite result, the empty set included:
\code{sqrt([1, DBL\_MAX])}, \code{nth\_root([DBL\_MAX], 3)} and the roots of
order 4 and more, \code{sqrt\_rel()}, \code{nth\_root\_rel()},
\code{sin([-1, DBL\_MAX])} and \code{cos()} (the width of the interval
overflowed), \code{tan([DBL\_MAX])}, \code{div\_rel()}, \code{atan\_rel()},
the reverse functions of \code{atan2} and of the power, \code{atanh\_rel()}
and, with Visual C++, \code{atan([DBL\_MAX])}. A quotient, a product or a
power that would overflow is replaced by what the operation gives, the
infinity or the largest double, without being computed; the roots of an
argument above $2^{1000}$ are proved on the argument and the root scaled by
powers of 2; \code{sin()} and \code{cos()} compare the half of the width
with $\pi$. A probe of 131 functions and forms on operands next to the limits
of the format found 34 that raised them for a finite result; none does with
the SSE2 intervals, the FPU intervals and the AVX-512 path. The bounds are the
same, bit for bit, up to the sign of a zero bound of \code{div\_rel()} with
the SSE2 intervals.

\medskip
The modes that flush subnormals to zero are cleared by each operation that
computes bounds.\newinvfive{} On x86 the flush-to-zero and denormals-are-zero
modes of MXCSR, on ARM the mode FZ of FPCR (FPSCR), make every subnormal
result a zero and read every subnormal operand as a zero. With one of them
set, the bounds of every operation with a subnormal operand or result miss
the exact result whatever the rounding direction: \code{[1e-300] * [1e-20]}
is \code{[0, 0]}. A program linked with \option{-Ofast}, \option{-ffast-math}
or \option{-funsafe-math-optimizations} has them (\file{crtfastmath.o}), and
so has one that loads a plug-in built so. Each operation that computes
bounds computes $1 + (2^{-1060} + 0)$ (x86-64, and ARM with GCC and Clang),
which shows the modes as well as the direction, and clears them with the
direction; a 32-bit x86 processor and the NEON intervals of 64-bit ARM
processors read the control registers instead, which hold both. With
\code{GAOL\_PRESERVE\_ROUNDING} they are set back as they were, and the
result, maxima and minima of bounds included, is made before. The check comes
before the operation reads, compares or copies a bound, and the compiler is
kept from moving a comparison before it: in the earlier versions of GAOL v5,
with GCC 13 at \option{-O3}, \code{asinpi([100*2\textasciicircum{}-1074])} was
above the exact value, and \code{sqrt([-1e-310, 4])} and \code{atan2} gave the
empty set. At the link, \file{gaol.pc} and
\code{gaol::gaol} give \option{-mno-daz-ftz} to the code using GAOL where the
compiler accepts it (GCC 13 and later on x86, from 11.4 and 12.4 in the series
11 and 12, and Clang 20 and 21, not Clang 18), which keeps \file{crtfastmath.o} out of the program; the check
of the operations is the only defence against a plug-in, and the only one on
the other compilers; on POWER, s390x, RISC-V and MIPS neither defence exists. The cost
is that of the second addition of the check: \code{x * y}, \code{sqrt} and
\code{pow(x, 3)} take 0.6 to 0.7 ns more than with the direction alone (Intel
Xeon Cascade Lake, Clang 18: 4.0 ns rather than 3.4, 11.0 rather than 10.3,
13.7 rather than 13.0). The check of the rounding direction alone,
$1 + 2^{-60}$, is the one left where GAOL cannot clear the modes (POWER,
s390x, RISC-V, MIPS, and the FPU intervals of Visual C++ for ARM)
(\pxref{sec:ftz}).

\medskip
The functions that only read, compare or copy bounds give right results under
these modes. With denormals-are-zero a comparison reads a subnormal as a
zero, and the functions that make no check gave wrong results:
\code{interval(2\textasciicircum{}-1073, 2\textasciicircum{}-1074)} kept its bounds in the wrong order,
\code{abs([-1e-309, -1e-310])} stayed negative, \code{floor([-1e-310])} was
\code{[-0]}, the midpoint of $[2, 4] \cdot 2^{-1074}$ was 0 and
\code{split()} made an empty half of it. The constructor from two bounds, the
relations, \code{\&}, \code{|}, \code{max}, \code{min}, \code{abs},
\code{sign}, the roundings to an integer, \code{mig}, \code{mag},
\code{is\_canonical}, \code{feven}, \code{nb\_fp\_numbers}, and \code{inf}
and \code{sup} of \code{gaol\_ieee1788} now compare the bounds as doubles
first and, where these are equal, by their bits, which no mode changes, and
round a subnormal bound from its bits. \code{bisect(ratio)} and
\code{is\_bisectable()} compare the bounds by their bits too, the cut point
included: $[3 \cdot 2^{-1074}, 5 \cdot 2^{-1074}]$ was not bisectable. The
AVX-512 path reads its operands once the modes are cleared, and makes its
results before they are set back: Clang 18 computed \code{x - y} after, and
\code{[3e-308] - [2.9e-308]} was \code{[0, 0]} (\pxref{sec:avx512}).

\subsection{The constants and the static objects}
\label{sec:changes-constants}
\cindex{constants!differences with GAOL~4}%

The constants of GAOL are initialized when compiling, so that the intervals
computed in the initialization of the static objects of a program are
right.\newinvfive{} They were computed by the dynamic initialization of
GAOL's files, which, with the static library of the CMake build, comes after
the static objects of the program, except with MinGW-w64: an interval
computed there found them 0, and \code{interval::pi()} was $[-0, 0]$,
$[1, 2]/[-1, 1]$ was $[-0, 0]$ rather than $[-\infty, +\infty]$ and
\code{sin([0, 4])} was $[-0, 0]$, without any warning (25 of the 45 checks of
the test of the static initialization failed with the SSE2 intervals, 24 with
the others). The doubles around $\pi$ and $\pi/2$, \code{gaol\_pi\_dn},
\code{gaol\_pi\_up}, \code{gaol\_half\_pi\_dn} and \code{gaol\_half\_pi\_up},
are written in decimal, exactly, rather than read from the unions that give
their bits, which is not a constant expression; they are in \code{gaol\_core}
(GAOL~4 had them in \code{gaol} as \code{pi\_dn}\ldots), and \code{gaol} and
\code{gaol\_ieee1788} name them too. \code{interval::one()}, \code{pi()},
\code{universe()} and the other constants build their interval from these
doubles, which the compiler folds, rather than copy the static intervals
\code{interval::cst\_one}, \code{cst\_pi}\ldots, which are gone
(\pxref{sec:interval-constants}). The masks of the SSE2 intervals are written
as vectors, which the compilers put in read-only memory (GCC and Clang) or in
initialized data (Visual C++), rather than computed with \code{\_mm\_set\_pd()}
and \code{\_mm\_cmpeq\_pd()}.

\medskip
\code{GAOL\_INFINITY} is \code{std::numeric\_limits<double>::infinity()}, a
constant of the compiler, where GAOL~4 took the \code{HUGE\_VAL} of the C
library.\newinvfive{} The UCRT of Windows (SDK 10.0.26100) writes
\code{HUGE\_VAL} as \code{((double)(float)1e+300)}: Visual C++ makes the
conversion when compiling, but clang-cl under \option{/fp:strict} makes it
when the program runs, in the rounding direction of the moment, and got
\code{FLT\_MAX} when that direction was downward or toward zero, with the
overflow flag. \code{interval()} was then $[-\mathrm{FLT\_MAX},
\mathrm{FLT\_MAX}]$, \code{interval(1e300)} was empty and $[1]/[-1, 1]$ did
not contain $10^{300}$ (\pxref{sec:floating-point-constants}).

\subsection{The kernels: SSE2, NEON, AVX-512 and the FPU}
\label{sec:changes-kernels}
\cindex{kernels!differences with GAOL~4}%

On 64-bit ARM processors the intervals are computed with NEON, with every
compiler of them: Visual C++, clang-cl and the Clang of MinGW-w64 and MSYS2 as
well as GCC, Clang and AppleClang under Linux and macOS.\newinvfive{} GAOL~4
computed the intervals of every ARM processor with the FPU, and 32-bit ARM,
whose NEON has no vector of doubles, still does. The two bounds are the two
lanes of a \code{float64x2\_t}, as the SSE2 intervals keep them in a
\code{\_\_m128d}, and \code{+}, \code{-}, \code{*}, \code{/}, \code{sqr()}
and \code{sqrt()} are kernels that take and return them in registers. On an
Apple M2 (AppleClang 21.0), \code{x + y} takes 0.90 ns rather than 1.15,
\code{x * y} 1.52 ns rather than 3.26, \code{x / y} 2.08 ns rather than 3.22,
\code{sqrt(x)} 2.04 ns rather than 5.17 and a line of 13 operations 16.6 ns
rather than 50.6, the second figures being those of the FPU intervals. The
bounds are the same, but for the sign of a zero bound, which the maximum of
64-bit ARM, ranking $+0$ above $-0$, can set otherwise: the lower bound of
$[-1, 0] \cdot [-1, 0]$ is $-0$ with NEON. Each operation checks the rounding
direction with one \code{mrs} (0.8 ns on the M2), which raises no exception
flag; with \code{GAOL\_PRESERVE\_ROUNDING}, an operation takes 27 ns rather
than 37 (\pxref{sec:neon}).

\medskip
The AVX-512 path carries the rounding direction in the
instruction.\newinvfive{} \code{x + y}, \code{x - y}, \code{x * y},
\code{x / y}, \code{sqrt(x)}, \code{sqr(x)}, the operators of an interval and
a double, and the integer powers beyond 1 and $-1$ are computed with the
rounding that the AVX-512 instructions carry in themselves, on x86-64 where
the compiler takes them, never GCC for Windows, which does not align the
stack for the AVX-512 registers (GCC bug 54412). They set no rounding
direction and raise no exception flag. The bounds are those of the SSE2
intervals bit for bit, but for the sign of a zero bound of \code{x * d}
where the zero bounds of the interval differ in sign. The three builds take
the path where \code{GAOL\_PRESERVE\_ROUNDING} is on and the machine building
runs the instructions, where it computes 2.9 times faster at the median on an
Intel i7-1185G7, and leave it out without the option, where the operations take
20\,\% longer;
the library then needs a processor with them, and the initialization stops a
program run on another one with a message (\pxref{sec:avx512}). Visual C++
and clang-cl for x64 take the SSE2 intervals and the AVX-512 path, as GCC and
Clang do, the kernels being \code{\_\_vectorcall} there so that their
\code{\_\_m128d} arguments go in registers.

\medskip
\code{*}, \code{/} and \code{sqr()} of the SSE2 intervals compute with no
branch on the signs of the bounds, and the basic operations hand the bounds
over in registers.\newinvfive{} GAOL~4 chose the product, the quotient and
the square by \code{switch}es and nested branches on the signs of the bounds,
and a branch predicted wrong costs more than the operation it chooses; its
operators read their operands and wrote their result through memory.
\code{x * y} is now the larger of four products, each rounded upward, in the
two lanes of an SSE2 register; \code{x / y} one division of the two lanes by
the denominators that masks drawn from the signs choose; \code{sqr(x)} two
products and a mask. Only an empty operand, an infinite bound or a divisor
holding 0 takes a branch. On an Intel i7-1185G7 with Clang 21, \code{x * y}
takes 2.3 ns rather than 10.0 with GAOL 4.2.3, \code{x / y} 2.5 ns rather
than 7.5, and Shekel 5 108 ns rather than 189 (\pxref{sec:comparison}). The
sign of a zero bound of a product or of a quotient may differ from GAOL~4's,
the value never: \code{maxpd} chooses between a $+0$ and a $-0$ that compare
equal, where GAOL~4 chose one product by the signs; $[-1, 0] \cdot [-1, 0]$
has the lower bound $+0$, where GAOL~4 had $-0$. The FPU intervals multiply,
divide and square without a branch too: GAOL~4 took one of nine cases of the
signs for a product, and went down a chain of up to eight tests for a
quotient. On an Apple M2 (AppleClang 21, \code{GAOL\_SIMD=OFF}), over the
random intervals of the benchmark, \code{x * y} takes 3.5 ns rather than
13.9, \code{x / y} 3.2 rather than 8.5, \code{sqr(x)} 1.8 rather than 6.9,
and Shekel 5 138 ns rather than 298. A line of products of operands of one
sign, which a processor predicts, takes more time (seven \code{*} in a row:
35.0 ns rather than 23.3, and 1.7 times longer with GCC 13 on x86-64 in the
continuous integration), and with \code{GAOL\_PRESERVE\_ROUNDING} the kernels
lose 7\,\% on operands of one sign (\code{x * y} 42.4 ns rather than 39.7) and
gain 11\,\% on operands of random signs (42.4 rather than 47.5). The sign of a zero
bound of such a product is not specified: GCC 13 for POWER9 compiles the
maximum as an \code{xsmaxcdp}, which returns its first operand. \code{x * y}
forms every corner product, so that a subnormal corner that is no bound of
the result raises \code{FE\_UNDERFLOW}, and a corner that overflows
\code{FE\_OVERFLOW} (\pxref{sec:kernels}).

\medskip
Several operations of the kernels were corrected or sped up.\newinvfive{}
\code{x -= x}, \code{x /= x} and \code{x \%= x} enclose their results in the
FPU intervals: GAOL~4 wrote one bound before reading it, and \code{x -= x}
was $[a - b, 2b - a]$ for $x = [a, b]$, \code{[0.1] /= [0.1]} was empty, and
normalizing a row by its pivot, \code{row[j] /= row[i]} down to $j = i$, did
it. A zero bound times an infinite bound counts as 0 in the products of the
FPU intervals, as in those of the SSE2 intervals: $[0, 1] \cdot [1, +\infty]$
is $[0, +\infty]$ in every build, where the FPU intervals gave the empty set.
The operators of an SSE2 interval with a double compute by the sign of the
double rather than through \code{interval(d)}: on an Intel i7-1185G7 (GCC
9.4, CMake Release), \code{x += d} takes 3.9 ns rather than 13.6,
\code{x *= d} 6.3 ns rather than 25.2 and \code{x /= d} 4.1 ns rather than
18.0, giving the same sets. \code{max(x, y)} and \code{min(x, y)} of the SSE2
and NEON intervals take both bounds of the result from one comparison of the
stored bounds: on an Intel i7-1185G7, \code{max(x, y)} takes 2.3 ns rather
than 8.5 with GCC 15.2 and 2.0 ns rather than 5.6 with Clang 21; on 64-bit
ARM they are 1.5 to 2.4 times faster (a Neoverse N2, an Apple M2). The
constructors from a register, \code{interval(const \_\_m128d\&)} and
\code{interval(const float64x2\_t\&)}, give the empty set when a lane is
$-\infty$, which is $l = +\infty$ or $u = -\infty$, the lower bound being
stored negated; both are explicit now.
\code{minimum()} and \code{maximum()} of two zeros give $-0$ and $+0$ with
every compiler: Visual C++ compiled \code{(b <= a) ? b : a} as \code{minsd},
which gives its second operand for two zeros.

\medskip
Square roots are the tightest intervals, as IEEE 1788-2015 requires of its
basic operations: the lower bound is the square root rounded upward when its
square is the argument, the double below it otherwise.\newinvfive{} GAOL~4
computed it as $x/\sqrt{x}$ rounded downward, one double below the tightest
bound for half of the doubles; \code{sqrt()} takes the same time as before
(9.0 ns on an Intel i7-1185G7, GCC 9.4), and is bounded whatever the rounding
of the C library's \code{sqrt}, which Visual C++ for 32-bit x86 rounds to
nearest in every rounding direction. \code{sqrt()} is no longer slowed down
by GCC on Intel processors: GAOL took the smallest subnormal from
\code{std::numeric\_limits<double>::denorm\_min()}, a \code{long double}
constant whose conversion GCC does not fold with \option{-frounding-math},
and the x87 unit converted it with a microcode assist at each use, about
100 ns on an i7-1185G7. \code{sqrt()} took 124 ns rather than 7.8, and
\code{nth\_root(x, 2)} and \code{nth\_root(x, -2)} 113 and 118 ns rather than
6.8 and 10.2 (GCC 15.2, Intel i7-1185G7); the constant is a literal now, and
the bounds are the same, bit for bit. \option{-ffloat-store} is added only
where doubles are still computed on the x87 unit: CMake gave it to GCC on
every target, and GCC then stored every double variable in memory. On an
Intel i7-1185G7 (GCC 9.4, CMake Release), \code{x + y} takes 3.0 ns rather
than 9.1, \code{x * y} 4.9 ns rather than 18.3, \code{sqrt(x)} 8.3 ns rather
than 41.3, and \code{log(x)} 73 ns rather than 110.

\medskip
The SSE2 intervals are computed on 32-bit Windows with GCC, Clang, Visual C++
and clang-cl.\newinvfive{} An SSE2 interval holds a \code{\_\_m128d} and has
to be aligned on 16 bytes, and the C runtime of 32-bit Windows aligns the
memory of \code{malloc()} and \code{new} on 8: a \code{std::vector} of
intervals crashed on its first \code{movaps} (MinGW-w64 14.2 and 15.2 for
x86). GCC and Clang compile the code using GAOL with
\option{-faligned-new=8}, which the three builds give to the library and to
the code that uses it, and \file{gaol/gaol\_config.h} stops the compilation of
the SSE2 intervals on 32-bit Windows without it. Visual C++ and clang-cl have
no such flag: their SSE2 intervals keep the two bounds as two doubles of
alignment 8, loaded with \code{movupd}, so that a program gives no flag. The
memory of the SSE2 intervals is released with the function matching the one
that allocated it (\code{GAOL\_MEMFREE()}): \code{\_aligned\_malloc()} and
\code{\_aligned\_free()} on Windows, so that the SSE2 intervals can be used
with MinGW-w64 and MSYS2 on x64. The placement \code{delete} of
\code{interval}, called when a constructor throws, freed the caller's memory;
it now leaves it alone.

\subsection[The constructors, the empty set and the interval functions]{The constructors, the empty set and the functions of intervals}
\label{sec:changes-intervals}
\cindex{constructors!differences with GAOL~4}%
\cindex{empty set!differences with GAOL~4}%

The constructors give the empty set where their arguments are not an
interval, as IBEX does: \code{interval(+oo)}, \code{interval(-oo)},
\code{interval(+oo, +oo)}, \code{interval(-oo, -oo)}, \code{interval(2, 1)}
and bounds that are NaN are all the empty set, and assigning an infinity to an
interval empties it.\newinvfive{} IEEE 1788-2015 has no interval
$[+\infty, +\infty]$ nor $[-\infty, -\infty]$: its constructor fails there,
and gives the empty set (10.5.8, 12.12.7). GAOL~4 built them as intervals of
one infinite point, whose midpoint, \code{realmax}, was outside them. The
intervals with one infinite bound are unchanged. On an Intel i7-1185G7 (GCC
9.4, CMake Release), \code{interval(a, b)} takes 1.28 ns rather than 1.24 ns:
its bounds are set in a register and behind a branch, since a pair written to
memory or conditional moves cost several nanoseconds per construction
(\pxref{sec:constructors}). The constructor from one \code{double},
\code{interval(double)}, is explicit, and so are \code{expression(double)}
and \code{expression(const interval\&)}: \code{interval x = 0.5;},
\code{return 0.5;} in a function returning an interval, \code{min(x, 0.5)}
and \code{interval r = sqrt(2.0);} do not compile, where GAOL~4 took the
double for a point interval without saying so, the \code{sqrt(2.0)} of the C
library among them. The assignment \code{x = d}, \code{x \&= d} and
\code{x |= d}, the relations between an interval and a \code{double}, the
arithmetic operators of an interval and a \code{double} and
\code{set\_contains(d)} have overloads of their own, which take
\code{interval(d)} as the conversion did: \code{x = 0.5} and \code{x < 0.0}
still compile, and give the same results. \code{interval(double, double)} is
not explicit.

\medskip
Integers are taken as the integers they are.\newinvfive{} An integer given
where GAOL takes a \code{double} was converted to a \code{double}, in the
rounding direction in effect, which an integer beyond $2^{53}$ is not:
\code{interval(9007199254740993LL)}, $2^{53} + 1$, was $[2^{53}, 2^{53}]$ or
$[2^{53} + 2, 2^{53} + 2]$, neither containing it. An integer of any type
but \code{bool} is now taken as the integer it is by the constructors, the
assignments, the arithmetic operators, the relations, \code{set\_contains()},
\code{expression(n)}, and \code{numsToInterval()} and \code{isMember()} of
\code{gaol\_ieee1788}: \code{interval(n)} is the tightest interval
containing $n$, the two doubles around it where $n$ is no double. The
arithmetic operators enclose the exact result, but not always with the
tightest bounds, for an $n$ that is no double: $[1, 2] + (2^{53} + 1)$ is
$[2^{53}, 2^{53} + 4]$. The exponents of \code{pow} and \code{pown}, and the
orders of \code{nth\_root}, \code{rootn} and \code{pownRev}, take an integer
of any type too: GAOL~4 had \code{pow(x, int)} only, so that
\code{pow(x, 3000000000u)} was \code{pow(x, -1294967296)}, and converted the
order of \code{nth\_root} to an \code{unsigned}, so that
\code{nth\_root(e, -2)} of an expression was the order 4294967294.

\medskip
The constructors from strings are gone.\newinvfive{} \code{interval(const
char*)} took any \code{const char*} and converted it implicitly:
\code{interval(0)} was ambiguous, 0 being a null pointer as well as a number,
and \code{x + nullptr} compiled, and crashed. \code{gaol::textToInterval()}
reads the string instead, with the same names of functions, and throws
\code{std::invalid\_argument} for a string that is no interval, and for a
null pointer, through an overload for \code{const char*}, where the overload
for \code{std::string} would build a \code{std::string} from the pointer,
which the standard leaves undefined. \code{interval(const char*, const
char*)} kept the bounds $[2, 1]$ of \code{interval("2", "1")}, which
\code{is\_empty()} called empty but whose sum with $[0, 1]$ was $[2, 2]$;
\code{textToInterval("0.1", "0.3")} takes the left bound of the first string
and the right bound of the second, and gives the empty set for bounds in the
wrong order. A result thrown away gives a warning: \code{sqrt(x);} alone
leaves \code{x} as it was, and the functions whose result is all they do
carry \code{GAOL\_NODISCARD}, \code{[[nodiscard]]} in C++17 and later and
the attribute of the compiler before; a cast to void silences it, and a
program that does not want it defines \code{GAOL\_NODISCARD} empty
(\pxref{sec:nodiscard}).

\medskip
The functions of intervals change on the empty set and on their
rounding.\newinvfive{} The intersection of disjoint intervals is the empty
set of \code{interval::emptyset()}, whose bounds are NaN: GAOL~4 kept the
largest left bound and the smallest right bound, $[3, 2]$ for
$[1, 2] \cap [3, 4]$, which \code{is\_empty()} took for empty but the
operations on the bounds did not, and $([1, 2] \cap [3, 4]) + [0, 1]$ was
$[3, 3]$. \code{width()} of the empty set is NaN, as \ieee{wid} of IEEE
1788-2015, rather than $-1$, and \code{chi()} of the empty set is the
positive NaN, printed \code{nan} with every compiler. \code{hausdorff()}
returns the tightest upper bound of the distance, and 0 for equal bounds,
infinite ones included: GAOL~4 computed \code{fabs(a - c)} in the rounding
direction of the caller and returned $+\infty$ as soon as a bound was
infinite, so that \code{hausdorff([1, +oo], [1, +oo])} was $+\infty$ and a
fixed-point loop on a box with an unbounded side never stopped. An infinite
bound in one interval only still puts them at $+\infty$, as in GAOL~4
(\code{hausdorff([1, +oo], [1, 2])}), and the function then returns before the
rounding direction is checked, 2.5 to 4.4 times as fast as
through the computation of the distance (Intel i7-1185G7, GCC 9.4 and Clang
18). \code{nb\_fp\_numbers()} numbers the doubles by the bits of the
magnitudes, $-0$ and $+0$ being one number: \code{nb\_fp\_numbers(-0.0, 1.0)}
was 13830554455654793217, and is 4607182418800017409, the number of the
doubles from 0 to 1. \code{mid()} returns the tightest interval enclosing the
midpoint, computed as \code{midpoint()} computes it with the rounding
outward: it computed $a/2 + b/2$ rounded upward and subtracted the half of
$b$ for its lower bound, which is not exact when $|b| < 2^{-1021}$, so that
\code{mid([0, 2\textasciicircum{}-1074])} was $[2^{-1074}, 2^{-1074}]$ and 45150 of the
180901 intervals whose bounds are multiples of $2^{-1074}$ from $-300$ to
$300$ times it did not enclose their midpoints. \code{rad()} and
\code{mid\_rad()}, \ieee{rad} and \ieee{midRad} of the standard, are new
(Chapter~\ref{chap:functions}).

\medskip
The tools that algorithms of intervals write for themselves are
new:\newinvfive{} \code{width\_enclosure()}, \code{inflate(r)},\hfill\break
\code{interval::centered\_form(m, r)}, \code{bisect(ratio)}, \code{is\_bisectable()},
and \code{hull(a, b)} and \code{intersect(a, b)}, the functions of
\code{a | b} and \code{a \& b}. The first three compute with the operations
of intervals, each bound rounded outward: \code{interval::centered\_form(m, r)}
encloses $[m - r, m + r]$, which \code{interval(m - r, m + r)} does not, both
of its bounds being computed upward; \code{width\_enclosure()} is the
enclosure of the exact width that \code{width()} bounds from above.
\code{bisect(ratio)} cuts at $l(1 - \mathit{ratio}) + r \cdot \mathit{ratio}$,
which does not overflow (for a bound of $2^{1023}$ or more in magnitude it is
computed from the halves of the bounds) and lies strictly between the bounds when
\code{is\_bisectable()}, and throws \code{std::domain\_error} for a ratio
outside $(0, 1)$ or a NaN; \code{is\_bisectable()} is true when a double
lies strictly between the bounds (\pxref{sec:tools}).

\subsection{The powers and the roots}
\label{sec:changes-powers}
\cindex{powers!differences with GAOL~4}%

The powers with a real exponent were corrected, after the fix of
Codac.\newinvfive{} \code{pow(I, e)} for a floating-point \var{e} called
\code{pow(I, int)}, which truncated the exponent: \code{pow([4], 0.5)}
returned $[1]$. It now computes an integer \var{e} with \code{pow(I, int)}
and any other \var{e} with \code{pow(I, J)}. \code{pow(I, J)} computed the
powers of the negative part of \var{I} on its magnitude: \code{pow([-4,-1],
[0.5])} returned $[-1, 2]$. It is now hybrid: a degenerate integer exponent
$[n]$ takes \code{pow(I, int)}, the \ieee{pown} of the standard, which keeps
the negative part of \var{I} and gives $0^0 = 1$; any other exponent takes
the \ieee{pow} of the standard, on the part of \var{I} in $[0, +\infty]$,
where $0^y$ is 0 for $y > 0$ and has no value for $y \leq 0$:
\code{pow([0], [0.5])} was $[0, 4.9 \cdot 10^{-324}]$ and
\code{pow([0], [-0.5])} was $[\mathrm{MAX}, +\infty]$
(\pxref{sec:pow-namespaces}).

\medskip
\code{pow(I, J)} takes the tightest bounds at the corners of the
box.\newinvfive{} $x^y$ increases with $y$ for $x > 1$ and decreases for
$x < 1$, increases with $x$ for $y > 0$ and decreases for $y < 0$: over a box
of bases above 0 its extrema are at corners, and the pow of CORE-MATH,
correctly rounded upward, is taken there, as the upper bound, and, as the
lower bound, the double below it, or $x^y$ itself where that is a double,
which is proved from the bits of $x$ and $y$: \code{pow([4], 0.5)} is $[2]$,
and \code{pow([3, 4], [2, 3])} is $[9, 64]$. At a corner where $x$ is 0 or
$+\infty$, or $y$ infinite, the bound is the limit of $x^y$, and CORE-MATH
is not called: \code{pow([4, +oo], 0.5)} is $[2, +\infty]$. GAOL~4 computed
$\exp(J \log I)$, which multiplied the relative width of $\log I$, a few
$2^{-52}$, by $|y \log x|$: \code{pow([2], [1023.5])} was 1425 doubles below
the exact value and 748 above, and \code{pow([0, 2\textasciicircum{}-1000], [-oo, -1])}
had a lower bound 617 doubles below $2^{1000}$. Over the million powers of
the benchmark of \file{doc/compare}, its intervals are no wider than those of
libieeep1788, which computes each bound with MPFR, where $\exp(J \log I)$
made them $5.9 \cdot 10^{-14}$ wider relatively. With Clang 18 on an Intel
i7-1185G7, over 100\,000 boxes, \code{gaol::pow(x, J)} with $x = [1/b, b]$
takes 84 ns (GAOL~4: 138 ns), with $x = [a, +\infty]$ 55 ns (GAOL~4: 99 ns)
and with $J = [a, +\infty]$ 54 ns (GAOL~4: 126 ns).

\medskip
Integer powers are the tightest bounds, of either sign.\newinvfive{}
\code{pow(x, n)} of \code{gaol}, \code{pown(x, n)} of \code{gaol\_ieee1788}
and the integer exponents of \code{pow(x, y)} and \code{pow(x, p)} give the
tightest bounds of $x^n$, or one double beyond where $x^n$ is within twice
the error bound of a double ($2^{-98}$ relatively at $n = 3$ and about
$2^{-62}$ at the most; at $n = 3$, one bound in $2^{46}$ is one double beyond
for an $x$ at random), and exact where $x^n$ is a double: \code{pow([2, 3],
3)} is $[8, 27]$. GAOL~4 rounded each product of its binary exponentiation
outward and was up to $n + 1$ doubles from the tightest bounds, 2, 3, 5, 6
and 8 doubles for $n = 3$ to 7, and took $x^{-n}$ for $1/x^n$, two roundings.
Both bounds are computed at once, in the two lanes of an SSE2 or a NEON
register with the FMA instructions, or in a pair of doubles otherwise: each
product is rounded, \code{fma} gives its rest exactly, and the power is
carried as a sum of two doubles whose error is bounded beforehand, so that
each bound is rounded once. A corner whose powers leave
$[2^{-968}, 2^{960}]$ is computed alone, scaled by powers of two: GAOL~4's
$1/x^n$ overflowed before the inversion, and \code{pow([10], -400)} was
$[0, 5.6 \cdot 10^{-309}]$ rather than $[0, 2^{-1074}]$. \code{pow(x, 3)}
takes 5.1 ns against 9.1 with GAOL 4.2.3 (Intel i7-1185G7, Clang 21). The
exponent 2 is \code{sqr(x)} in every form, the squares of the bounds rounded
outward: in the earlier versions of GAOL v5, with the exact products,
\code{pow(x, 2)} took 5.3 ns with the NEON intervals and 5.2 with the FPU intervals, and it takes 1.6 and 1.9 now (Apple
M2, AppleClang 21); the 6 million calls of \code{pow(x, 2)} of the
Goldstein-Price example take 50 ms rather than 77 with the NEON intervals.
The integer powers beyond the ints are computed from the pow of CORE-MATH at
the corners of $x$, the sign from the parity of $n$: the tightest where $n$
is a double, and within three doubles of it beyond $2^{53}$. \code{pow(x, 2)} in all its
forms, the powers with a real or an interval exponent (an empty exponent
included), with a \code{long long} beyond the \code{unsigned} ints, and
\code{gaol\_ieee1788::pow} return the empty set for an empty
operand before they check the rounding direction, as \code{sqr} does, and leave
the direction and the flush modes as the program had them; the other integer
powers, of an \code{int}, an \code{unsigned} or a \code{long long} within their
range, check it first and leave it upward (\pxref{sec:powers}).

\medskip
\code{nth\_root(I, n)} is the \ieee{rootn} of IEEE 1788-2015: for an odd
$n$, it is defined on the whole real line, the root of a negative number
being the opposite of the root of its magnitude, and \code{nth\_root([-8, 27],
3)} encloses $[-2, 3]$; for an even $n$, it takes the roots of the part of
\var{I} in $[0, +\infty]$, as before.\newinvfive{} GAOL~4 took that part for
every $n$, and gave $[-2^{-1074}, 3.000000000000001]$, or the empty set for
\code{nth\_root([-8, -1], 3)}; the root of 0 is 0, rather than
$[-2^{-1074}, 2^{-1074}]$. The $n$th roots are proved with integer powers,
and are the tightest bounds or up to 2 doubles beyond (4 for \ieee{rootn}
with $q < 0$, the inverse of a root), for every $n$ and every double, the
subnormal ones included. GAOL~4 took the power of the mathematical library
with the exponent $1/n$ rounded, moved one double outward: the bounds were up
to 8 doubles from the tightest between $2^{-30}$ and $2^{30}$, and 234 over
all the doubles, and \code{nth\_root([27], 3)} was two doubles wide where it
is $[3]$ now. A bound $l$ is below the root when $l^n$, rounded upward, is at
most $x$, and $u$ above it when $u^n$, rounded downward, is at least $x$; the
search starts from the pow of CORE-MATH brought next to the root by a step of
Newton's method, and goes by steps that double, then by bisection. For an
$x$ below $2^{-1020}$ the proofs take scaled powers, the plain proof losing
bits where a power is subnormal: \code{nth\_root([2\textasciicircum{}-1074], 33)} was 1983
doubles below the tightest lower bound and 28785 above the tightest upper
bound. \code{nth\_root\_rel()} takes the roots of \code{nth\_root()}: it kept
a value within 23 doubles, now 2. A point interval takes 152 ns rather than
130, its only root being looked for once; an interval 194 ns rather than 136
(Intel i7-1185G7, GCC 9.4). \code{pownRev(c, x, p)} for $p \leq 0$, the last
reverse operation of the standard GAOL lacked, is \code{nth\_root\_rel(J, p,
I)}, which takes any integer $p$ now: a negative \code{int} was converted to
a large \code{unsigned int}, a root of order about $4 \cdot 10^9$. For
$p < 0$ the bounds are up to 4 doubles beyond the tightest ones.
\code{uipow()}, the \ieee{pown} for an unsigned exponent, is \code{pow(I,
n)} for an \code{unsigned} $n$ in \code{gaol}, exported with the SSE2
intervals too, where it was inline and did not link; it gives $[1]$ for an
exponent 0, and the empty set for an empty interval.

\subsection{The trigonometric, hyperbolic and other functions}
\label{sec:changes-functions}
\cindex{elementary functions!differences with GAOL~4}%

\code{sin()} is computed as \code{cos()} is: the bounds of the interval
divided by an enclosure of $\pi$, minus $1/2$, tell the pieces where the sine
is monotonic, and the sine of CORE-MATH is taken at the bounds.\newinvfive{}
GAOL~4 computed \code{cos(x - [pi/2])}, whose subtraction widened the
argument by about $2^{-52} \max(2, |x|)$: \code{sin([1e-10])} was
$4.4 \cdot 10^{-16}$ wide, billions of doubles, and \code{sin()} was not
accurate in the sense of IEEE 1788-2015. On an Intel i7-1185G7 (GCC 9.4),
\code{sin()} takes 121 ns rather than 130. \code{sin()}, \code{cos()} and
\code{tan()} are within one double of the tightest bounds at every
magnitude. The bounds divided by an enclosure of $\pi$ only tell on which
pieces the interval lies while the quotients are not within their rounding
errors of an integer: a bound within about $|x| \cdot 2^{-52}$ of an extremum
was taken as reaching it, \code{cos([2\textasciicircum{}60])} was $[-1, 1]$, and
\code{tan()} gave $[-\infty, +\infty]$ next to a pole and beyond $2^{52}$.
CORE-MATH being correctly rounded, and no double but 0 being a multiple of
$\pi/2$, the signs of its sine and cosine at the bounds are the exact ones,
and tell whether an extremum or a pole is within the interval; they are only
asked for next to an extremum or a pole, and at the large magnitudes. The
time is unchanged in the benchmark (113 ns for \code{sin()} and \code{cos()},
Intel i7-1185G7, GCC 9.4), 125 ns rather than 121 for intervals holding an
extremum; at magnitudes from $10^8$ to $10^{15}$, \code{cos()} of a narrow
interval takes 250 ns rather than 224, and \code{tan()} 260 rather than 217.
\code{tan()} is finite over an interval of doubles that holds no pole and
whose width, rounded upward, is at most the double below $\pi$: \code{tan([-M\_PI\_2, M\_PI\_2])} was
$[-\infty, +\infty]$, and is about $\pm 1.63 \cdot 10^{16}$
(Chapter~\ref{chap:accuracy}).

\medskip
\code{atan2(y, x)} is implemented, as the \ieee{atan2} of the standard,
defined on the plane but $(0, 0)$ with values in $(-\pi, \pi]$: GAOL~4
declared it, and its parser read it, but it raised
\code{unavailable\_feature\_error}.\newinvfive{} In each quadrant the angle
is monotonic in $y$ and in $x$, and the atan2 of CORE-MATH is taken at two
corners of the box; a box with points on the half-line $y = 0$, $x < 0$ and
points below it gives $[-\pi, \pi]$, and \code{atan2([0], [0])} is empty. It
takes 100 ns on an Intel i7-1185G7 (GCC 9.4) for boxes within a quadrant,
against 68 ns for \code{atan()}. \code{acos\_rel()}, \code{asin\_rel()} and
\code{atan\_rel()} enclose the multiples of $\pi$ of the pieces of the
preimage with $\pi$ in double-double, one rounding at the magnitude of the
result rather than two: \code{atan\_rel()} was up to 49 doubles from the
tightest bounds over random cases where it is now within 2, and
\code{acos\_rel()} up to 34. An \var{I} of a single double is decided at
every magnitude by whether $f(I)$ meets \var{J}, where a bound beyond
$2^{52}$ was kept as it was (\pxref{sec:relational-arithmetic}).

\medskip
The elementary functions are exact where their value is a double, and the
tightest where it is $\pi/4$, $\pi/2$ or $\pi$.\newinvfive{} $\exp(0) = 1$
and $\log(1) = 0$ exactly, the bounds moved one double outward giving
\code{exp([0])} and \code{log([1])} a width: \code{pow([1], [-oo, +oo])} was
$[0, +\infty]$ rather than $[1]$. \code{log()} takes 1.5 ns more (72.5 ns
rather than 71) and \code{pow(x, y)} 2 ns more (145 ns rather than 143) on
an Intel i7-1185G7 (GCC 9.4). \code{sin(0)}, \code{tan(0)}, \code{asin(0)},
\code{atan(0)}, \code{sinh(0)}, \code{tanh(0)}, \code{asinh(0)} and
\code{atanh(0)} are 0, \code{cos(0)} and \code{cosh(0)} 1, \code{acos(1)}
and \code{acosh(1)} 0, and \code{asin(}$\pm$\code{1)}, \code{acos(0)},
\code{acos(-1)}, \code{atan(}$\pm$\code{1)} and \code{atan(}$\pm$\code{oo)}
are the constants of GAOL: \code{acos([1, 3])} was
$[-2^{-1074}, 2^{-1074}]$, and \code{acosh([1])} with the libm was
$[-3 \cdot 2^{-1074}, 3 \cdot 2^{-1074}]$. The hyperbolic functions
\code{sinh}, \code{cosh}, \code{tanh}, \code{asinh}, \code{acosh} and
\code{atanh} are those of CORE-MATH, the tightest bounds on every system:
GAOL~4 took them from the libm of the system and moved their values one double
outward, while the libms of glibc 2.31, musl and MinGW-w64 are sometimes a
double further, and the \code{acosh} of the MinGW-w64 older than version 12 is
millions of doubles away next to 1. \code{cosh()} and \code{sinh()} are
$[\mathrm{MAX}, +\infty]$ from 711 on, and \code{tanh()} has the double below
1 for lower bound from 20 on; the functions take the same time as before,
within 1.5 ns (Intel i7-1185G7, GCC 9.4).

\medskip
The domains follow the standard.\newinvfive{} \code{log()} gives the empty
set for an interval holding no positive number, \code{log} being defined on
$(0, +\infty)$: GAOL~4 gave $[-\infty, -\mathrm{MAX}]$ for \code{log([-4,
0])} and \code{log([0])}. \code{atanh()} gives the empty set for an interval
holding no point of $(-1, 1)$: \code{atanh([1])} and \code{atanh([1, +oo])}
are empty, and the limits stay the bounds of an interval holding other
points, \code{atanh([0, 1])} being $[0, +\infty]$. \code{div\_rel()} rounds
the lower bound of the quotients downward where \var{K} is below 0 and
\var{J} straddles it, the one branch left without SSE code:
\code{div\_rel([-1], [-3, 1], [0, 10])} was one double above $1/3$, and the
numbers from $1/3$ to the double above it were lost, where \ieee{mulRev} has
to be valid. \code{invabs\_rel()} takes no magnitude from the negative part
of \var{J}: \code{invabs\_rel([-2, -1], [-10, 10])} was $[-2, 2]$ rather than
the empty set. The reverse functions of the power, \code{atan2}, \code{max},
\code{min}, \code{sign} and \code{floor} are new, their arguments in the
order of \code{div\_rel(K, J, I)}: \code{pow\_rel(K, J, I)},
\code{pow\_exponent\_rel(K, H, I)}, \code{atan2\_rel\_y(C, X, Y)},
\code{atan2\_rel\_x(C, Y, X)}, \code{max\_rel(Z, Y, X)},
\code{min\_rel(Z, Y, X)}, \code{sign\_rel(Z, X)} and \code{floor\_rel(Z,
X)}; \code{gaol\_ieee1788} names the first four \ieee{powRev1},
\ieee{powRev2}, \ieee{atan2Rev1} and \ieee{atan2Rev2}, and adds
\ieee{mulRevToPair}. All the reverse operations of the standard are
provided (Chapter~\ref{chap:arithmetic}).

\subsection{The relations}
\label{sec:changes-relations}
\cindex{relations!differences with GAOL~4}%

The comparisons with the empty set and the infinities follow IEEE 1788-2015
(Tables~10.3 and 10.4).\newinvfive{} \code{certainly\_leq()} and
\code{certainly\_le()}, \ieee{precedes} and \ieee{strictPrecedes}, and
\code{certainly\_geq()} and \code{certainly\_ge()}, are true when either
interval is empty: GAOL~4 gave false when the second one only was.
\code{set\_strictly\_contains()} and \code{set\_le()}, \ieee{interior}, take
an infinite bound as beyond the same infinite bound:
\code{interior(Entire, Entire)} and $[2, +\infty]$ interior to $[1, +\infty]$
were false, as \code{set\_le()} of the empty set in the empty set.

\medskip
The possibly relations, \code{certainly\_eq()}, \code{certainly\_neq()},
\code{==} and \code{!=} are removed, and so is the option that chose what the
relation symbols mean.\newinvfive{} \code{certainly\_neq()} was
\code{!certainly\_eq()}, which is ``possibly not equal'': $[3, 4]$ was
certainly not equal to $[3, 4]$, and so was \code{!=} with the certainly
relations, the default of every build. IEEE 1788-2015 has neither the
possibly relations nor a certain equality: its equality is \ieee{equal},
\code{set\_eq()}, and ``for all $x$ and $y$, $x \neq y$'' is
\ieee{disjoint}, \code{set\_disjoint()}; \code{possibly\_eq(y)} was
\code{!set\_disjoint(y)}. \code{<}, \code{<=}, \code{>} and \code{>=} are
the certainly relations in every build, \ieee{strictPrecedes},
\ieee{precedes} and their converses. The set relations of the option were
never taken, \cmd{configure} and meson defining \code{GAOL\_SET\_RELATIONS}
where the header read \code{GAOL\_SET\_RELATION}: they gave the possibly
relations (\pxref{sec:relational-symbols}).

\medskip
\code{gaol::lexicographic\_less} is a total order on the intervals, for the
containers and the algorithms of the standard library that compare with
\code{<} (\code{std::set}, \code{std::map}, \code{std::sort}).\newinvfive{}
\code{<} serves none of them: it is \ieee{strictPrecedes}, true as soon as
either interval is empty, and no strict weak ordering. A
\code{std::set<interval>} that holds $[1, 2]$ drops $[1, 3]$, and
\code{std::sort} of a vector that holds an empty interval can read past its
end. With the comparator, the empty set comes first, then the intervals by
their left bound and by their right bound, $-0$ and $+0$ being the same bound
(\pxref{sec:lexicographic-less}). \code{x.overlap(y)} and \code{gaol\_ieee1788::overlap(x, y)}, the
\ieee{overlap} of IEEE 1788-2015 (10.6.4), are new, with the type \code{overlapping\_state}: each returns one of
the 16 states of Table~10.7, which GAOL~4 had no function for. The state is
read from the bounds as integers in the order of the doubles, so that no
floating-point exception is raised, whatever the flush-to-zero modes. A point
interval on a bound of the other one is never \element{meets} nor
\element{metBy}, as the sentence of 10.6.4 has it
(Chapter~\ref{chap:relations}).

\subsection{Reading and writing intervals}
\label{sec:changes-io}
\cindex{input!differences with GAOL~4}%
\cindex{output!differences with GAOL~4}%

Numbers are read exactly, whatever the locale, the flush-to-zero modes and
the C library.\newinvfive{} GAOL~4 read them with \code{strtod()} rounding
downward and upward, and relied on the inexact flag: the C runtime of Windows
and musl on 64-bit ARM processors round to nearest in every direction and
raise no flag, so \code{interval("0.1")} did not enclose $1/10$ there. Each
number is now compared exactly with the doubles around it, whatever the
number of digits of its significand and of its exponent; the doubles around
the number are bracketed by steps doubling from the double \code{strtod()}
gives, then found by bisection on their bits, two comparisons when
\code{strtod()} gives the nearest double, and at most about 125 whatever it
gives. The time to read a number is quadratic in the number of its digits
(\pxref{sec:reading-intervals}), so that a text from elsewhere can block a
program.
\code{strtod()} reads the decimal point of the C locale: under a locale
writing a decimal comma, it stopped at the \code{.} of \code{0.1}, and
\code{interval("0.1")} was $[0, 0]$; the text is taken apart once, and a
number of 20000 characters, which took over 2 s under such a locale, takes
about 0.02 s. Reading raises no overflow or underflow exception: with glibc,
\code{strtod()} raises them for \code{1e400} or \code{1e-310}, and a program
that had enabled the exception died while it read; a number that may be
beyond the largest double or below $2^{-1022}$ is compared with them first,
and is not given to \code{strtod()}. With denormals-are-zero the reader took
every subnormal for 0, and \code{textToInterval("1e-310")} did not contain
$10^{-310}$; the doubles are taken apart by their bits, which no mode
changes (\pxref{sec:reading-intervals}).

\medskip
Intervals are written in decimal rounded outward, whatever the C
library.\newinvfive{} GAOL~4 set the rounding direction downward, then
upward, and let the C library write each bound: the C runtime of Windows
rounds the magnitude, and wrote $-2/3$ rounded downward \code{-0.6666}, and
musl wrote \code{1.235e+05} for the lower bound of 123456.789; read back,
such a text did not enclose the interval written. The magnitude of each
bound is now written rounding to nearest, compared exactly with the bound,
and its last digit is moved by one when it is on the wrong side; with glibc,
which rounds as asked, the text is the one the C library wrote: 1\,077\,546
bounds written in the general, scientific and fixed formats with 1 to 18 digits
were the same, character for character. A
point interval whose double the digits write exactly is written $[a]$, the
literal of the standard, and a zero of either sign $[0]$: GAOL~4 wrote every
point interval \code{<a, b>}, the text of its double rounded downward and
upward, which the reader refused for \code{interval(0.1)}, throwing
\code{input\_format\_error}, because it takes \code{<a, b>} for two numbers
that are the same double, so that reading back the text of a point interval
failed for about 61\,\% of 11\,788 point intervals tried; the reader still reads \code{<a, b>}, with the numbers
compared by their bits. The format \element{width} writes the midpoint and
the radius of the standard, \code{midpoint()} rounded to nearest and
\code{rad()} rounded upward, which is $+0$ for a point zero of either sign, and
\element{center} the midpoint: GAOL~4 wrote
$(l + r)/2$ and $(r - l)/2$, both rounded to nearest, so that $[1, 1 +
2^{-52}]$ was written \code{1 (+/- 1.110223024625157e-16)}, which stops short
of the upper bound, and $[0, 5 \cdot 10^{-324}]$ had a radius 0. The empty
set is \code{[empty]} in every format. The format \element{bounds} is
written in the C locale whatever the locale of the stream, so that the reader
reads the text again: GAOL~4 wrote \code{[0,25, 0,5]} under a comma locale.
The format \element{agreeing} keeps the digits of each bound: GAOL~4 dropped
from both bounds what followed the last character of the left one that is
not a zero, and $[1.25, 1.2567]$ was written \code{1.25\textasciitilde{}[, ]}, which says
nothing of the right bound; it is \code{1.25\textasciitilde{}[0, 67]} now, and $[1, 10]$
is written with its bounds, where it was \code{1\textasciitilde{}[., 0]}. The hexadecimal
format writes an interval literal, \code{[0x1.999999999999ap-4, ...]}, read
back bit for bit, where GAOL~4 wrote the sixteen hexadecimal digits of each
double, which no parser read; \code{exact\_string(I)} gives that text
without the global output format, and \code{gaol\_ieee1788::intervalToText()}
writes a literal with the digits of \code{interval::precision()}, whatever
the format and the locale. \code{operator<<} leaves the stream as it found
it: GAOL~4 set its precision and left it so, the doubles written afterwards
getting 16 digits, and wrote the \code{[} apart, so that
\code{std::setw(8)} padded it alone. The text is made in a character string,
taking 30 to 45\,\% less time than the \code{std::ostringstream}s of the earlier
versions of GAOL v5 (Clang 18, Release) (\pxref{sec:writing}).

\medskip
\code{operator>>} ends with the input, and skips blank lines.\newinvfive{}
GAOL~4 read the rest of the line where the previous value stopped, so that
\code{in >> d >> x} over two lines threw \code{input\_format\_error}, and so
did \code{while (in >> x)} over a file ending with an empty line: such a loop
always ended with an exception. Where no line is left, the reader sets the
\code{failbit} of the stream, leaves the interval as it was and throws
nothing, as the reading of a double does; a line that is no interval sets the
\code{failbit} too, empties the interval, and \code{std::invalid\_argument}
is thrown. A program that reads on after it calls \code{in.clear()} first.

\medskip
One grammar reads every string as an expression.\newinvfive{} GAOL~4 read a
string with two grammars: the numbers, the constants and the functions of
them built a tree, while the intervals written between brackets, and the
operations on them, were computed at once; Bison chose between the two at
each token (124 shift/reduce and 14 reduce/reduce conflicts), and once a
number had started the tree no interval could come, so that \code{1+[1,2]}
and \code{2*cos([0,1])} were refused while \code{[1,2]*2} was read. The
literals of intervals, the uncertain numbers, \code{empty} and \code{<a,b>}
are now leaves, computed when they are read, as are the calls of functions
and the binary operations; the grammar has no conflict, and an error found
in an action stops the reading at once. An uncertain number is still no
bound of a literal: \code{[5?1]} throws \code{std::invalid\_argument}, the
standard giving it as a string that is not an interval literal. The literals
of IEEE 1788-2015 are read, whatever the case of their letters: \code{[ ]},
\code{[entire]}, the bounds left out (\code{[1,]}, \code{[,]}),
\code{infinity}, the hexadecimal numbers and the uncertain form
(\code{3.56?1}, \code{-10??u}), as the tightest intervals of doubles
enclosing them; \code{[inf]} and \code{[inf, inf]} are the empty set, where
GAOL~4 gave $[\mathrm{MAX}, +\infty]$. Numbers in base 16, 8 and 2 are read,
\code{0x10}, \code{0o17}, \code{0b101} and \code{0b1.1p3}, a number that
begins with 0 being decimal. $[l, u]$ with $l$ above $u$ is the empty set
also where the bounds are numbers that are not doubles: GAOL~4 read
\code{[1.00000000000000001, 1]} as $[1]$. A bound of \code{<a, b>} has to be
a double, and a number that is not one, \code{pi} included, is refused. The
reader knows every function of GAOL on intervals, which GAOL~4 did not read
(\code{exp10}, \code{hypot}, \code{sqr}, \code{abs}, \code{fma}\ldots), and
\code{gaol\_ieee1788::textToInterval()} reads the names of the standard,
whose \code{pow} is the pow of the standard: in the earlier versions of GAOL v5
it read those of GAOL, and \code{pown} gave the empty set. Function names are read in any case of
letters, compared in ASCII and not by \code{std::tolower()}, whose result
depends on the locale. \code{pow(x, n)} read in a string is the \code{pow}
of \code{gaol} for an integer exponent too: the parser took $1/x^{|n|}$ for
$n < 0$, whose $x^{|n|}$ overflowed, and 1 for $n = 0$, which is no value for
an empty $x$ (\pxref{sec:input-format}).

\medskip
The reader is robust.\newinvfive{} The parser frees the nodes of the
expressions it reads: it kept a reference to each node, and the nodes of
every \code{interval("...")} were never freed; the nodes left by a syntax
error are freed too, and reading \code{nth\_root(8, 1.5)}, then the same
text plus 1, crashed, where both now throw \code{std::domain\_error}.
Several threads can read strings at once: the lexer of flex kept its buffer
in globals, and two threads each building an interval from a string crashed;
the lexer is reentrant and the parser pure, and four threads reading strings
at once, which crashed each time with the globals, show nothing under
ThreadSanitizer. Long expressions no longer overflow the stack: a sum of
100000 terms crashed the program, and is read in a time proportional to its
length; the stack of the parser, 10000 entries, bounds the nesting. A
\code{std::string} is read whole, where the lexer stopped at the first NUL
character. An allocation that fails throws \code{std::bad\_alloc}, where
flex ended the program with \code{exit(2)}. Built without exceptions (a setting of the developers), a
wrong call in a string fails the reading as a syntax error does, where an
earlier version of GAOL v5 ended the program inside the reader:
\code{gaol\_ieee1788::textToInterval()} gives the empty set,
\code{parse\_interval()} returns false, and \code{gaol::textToInterval()} and
\code{operator>{}>} report the error and abort. \code{parse\_interval(nullptr)} returns false, where the program
died of a segmentation fault. A newline is a space, where it fell to the
default rule of flex, which wrote it on the standard output. Expressions are
printed as written: a division with \code{/}, where GAOL~4 wrote \code{*},
and a negative number within the parentheses a power of it needs.

\subsection{The errors}
\label{sec:changes-errors}
\cindex{exceptions!differences with GAOL~4}%

GAOL throws \code{std::invalid\_argument} and \code{std::domain\_error}, and
\code{what()} says what went wrong.\newinvfive{} GAOL~4 had a class of
exceptions of its own, \code{gaol\_exception}, and
\code{input\_format\_error}, \code{invalid\_action\_error} and
\code{unavailable\_feature\_error} derived from it, in
\file{gaol/gaol\_exceptions.h}. They are gone, without aliases. A string
that is no interval, or a null pointer, throws \code{std::invalid\_argument};
an argument that a function does not take (\code{nb\_fp\_numbers()} with a
NaN, \code{bisect()} with a ratio outside $(0, 1)$, \code{nth\_root(8, 1.5)}
in a string) throws \code{std::domain\_error}; no operation of GAOL v5 throws
\code{unavailable\_feature\_error}, which only the \code{atan2} of GAOL~4
did. A handler of \code{gaol::input\_format\_error} becomes a handler of
\code{std::invalid\_argument}, one of \code{gaol::invalid\_action\_error} a
handler of \code{std::domain\_error}, and \code{std::logic\_error} catches
both. GAOL~4 left the \code{what()} of \code{std::exception}, the text
\code{std::exception} with libstdc++ and libc++ whatever the error; it is
\code{file:line: explanation} now, the file without its directories.
\code{file()}, \code{line()}, \code{explanation()} and \code{operator<<} of
the exceptions are gone with the classes. A null file name is ``(unknown
file)'' in the message of an exception, in \code{gaol\_warning()} and in
\code{gaol\_error()}: GAOL~4 wrote the null pointer to the stream, which is
undefined, and libstdc++ put the stream in the state \code{bad}, so that the
diagnostics written after it were lost (\pxref{sec:exceptions}).

\subsection{The namespaces, the headers and the interface}
\label{sec:changes-interface}
\cindex{namespaces!differences with GAOL~4}%
\cindex{headers!differences with GAOL~4}%

The type \code{interval}, GAOL's functions and its expressions are in
\code{gaol\_core}; \code{gaol} names them as GAOL~4 did, and
\code{gaol\_ieee1788} as IEEE 1788-2015 does.\newinvfive{} \code{gaol} names
them one by one, with using-declarations, rather than opening
\code{gaol\_core} whole: the names GAOL~4 had in \code{gaol} that GAOL v5
still has, the functions on intervals GAOL v5 adds, \code{restore\_rounding()}
and \code{exact\_string()}. The two namespaces differ mainly in the power,
the standard having two where GAOL~4 had one: the \code{pow} of \code{gaol}
takes the integer power for an integer exponent, a negative base included,
and the real power otherwise, while in \code{gaol\_ieee1788} \code{pow} is
the real power, on the part of $x$ in $[0, +\infty]$, and \code{pown} the
integer power. \code{pow} and \code{textToInterval}, whose meaning differs
between the two, are in each of them and not in \code{gaol\_core}, so that
each namespace finds its own only. \code{gaol} does not name the exceptions,
\code{NaN\_val}, \code{eval\_stack} nor \code{parse\_interval()}, which
GAOL~4 had in \code{gaol}, nor what GAOL uses for itself; a program names
them in \code{gaol\_core}. \code{f\_negate()}, \code{f\_negate\_simple()},
the macro \code{opposite()}, \code{get\_fpu\_cw()} and
\code{reset\_fpu\_cw()}, which GAOL v5 no longer calls, are gone.
\file{gaol/gaol.h} no longer opens \code{gaol}: a program adds \code{using
namespace gaol;} or \code{using namespace gaol\_ieee1788;}, one of the two,
\code{pow(x, y)} being ambiguous with both (\pxref{sec:namespaces}). The
\code{double} constants of \file{gaol/gaol\_port.h}, \code{two\_pi},
\code{pi}, \code{half\_pi}, \code{pi\_dn}, \code{ln2\_dn}\ldots, are no longer
declared for the code using GAOL: a program that opened \code{gaol} found
them among its own names, and a \code{pi} of its own was ambiguous; the
bounds of $\pi$ are those of \code{interval::pi()}, and those of $\ln 2$
those of \code{log(interval(2.0))}.

\medskip
The expressions were repaired.\newinvfive{} \code{pow(e, n)} on an
expression links: the library defined it with an \code{unsigned int}
exponent, where the header declares an \code{int}. \code{expression::operator/=}
links: it was declared and not defined. Empty expressions can be created in
several threads: they all point to one node, whose count of references, a
plain \code{unsigned int}, lost updates, and two threads each creating and
destroying a million empty expressions crashed each time; the references to
that node are no longer counted, and it lives from the initialization of GAOL
to its automatic cleanup. Other expressions still share the nodes of their
copies, whose counts are not atomic: such an expression is not copied or
destroyed in two threads at once.

\medskip
The headers GAOL installs are the ones a program includes.\newinvfive{} The
three builds install the same list: \file{gaol/gaol}, \file{gaol/gaol.h}, the
headers they include, and \file{gaol/parser/gaol\_expression.h} and
\file{gaol/parser/gaol\_expr\_eval.h} for the expressions; a program that
builds expressions includes \file{gaol/parser/gaol\_expression.h}, where it
included \file{gaol/gaol\_expression.h}. They install no header of GAOL's sources:
neither the header of the parser, which did not compile when included, nor
\file{gaol/gaol\_init\_cleanup.h}, both of which GAOL~4 installed, nor the
headers of \file{gaol/math/}, among them \file{gaol\_double\_op.h}, whose
functions of doubles \code{exp\_dn()}, \code{sin\_up()} and the thirty others
leave the interface. Each installed header, but the four of Section~\ref{sec:headers}, compiles
alone with
\option{-std=c++11} and the flags of \file{gaol.pc}, and without a warning
under \option{-Wall -Wextra} and \option{-Wold-style-cast -Wundef}: with
\option{-Wextra}, every assignment of intervals warned that the implicit
assignment was deprecated. The class \code{interval} is declared in
\file{gaol/core/gaol\_interval\_core.h}, and the headers of the SSE2, NEON
and FPU intervals are installed in \file{include/gaol/core/};
\file{gaol/gaol\_interval.h} stays and includes it, so that a program that
includes it, or \file{gaol/gaol}, is not touched by the move. Gone are the
timing of GAOL~4 (\file{gaol/gaol\_profile.h}, the class \code{timepiece}
and \code{get\_time()}, \code{elapsed\_time()}\ldots; a program measures the
time with the clocks of the C++ standard library),
\file{gaol/gaol\_parameters.h} and \file{gaol/gaol\_limits.h} (a program
includes \file{<limits>} itself), \file{gaol/gaol\_assert.h} and
\code{GAOL\_ASSERT} (a program writes \code{assert()}),
\file{gaol/gaol\_exact.h} with \code{get\_inexact()} and
\code{clear\_inexact()}, which the lexer called to know whether
\code{strtod()} had rounded, \file{gaol/gaol\_flags.h} and
\code{s\_nextafter()}, and \file{gaol/gaol\_allocator.h} with
\code{gaol::aligned\_allocator}, which neither GAOL, Codac nor IBEX used.
The \code{operator new} of the SSE2 intervals still aligns their memory.

\medskip
\code{\_\_GAOL\_PUBLIC\_\_} is \code{GAOL\_PUBLIC}, and is empty with Visual
C++ and clang-cl, with which GAOL is always a static library.\newinvfive{}
GAOL~4 declared its functions \code{\_\_declspec(dllimport)} there, so that
the code using GAOL had to define \code{\_\_GAOL\_PUBLIC\_\_} empty; no
header reads it any longer. No header of GAOL declares \code{string} and
\code{exception} in the global namespace, and the macros of GAOL's headers
start with \code{GAOL\_}: \code{INLINE}, \code{MEMALIGN}, \code{\_\_HI},
\code{WORDS\_BIGENDIAN}\ldots are \code{GAOL\_INLINE},
\code{GAOL\_MEMALIGN}, \code{GAOL\_HI}, \code{GAOL\_WORDS\_BIGENDIAN}\ldots
\code{\_XOPEN\_SOURCE} is no longer redefined to 600, and
\code{nb\_fp\_numbers()} returns an \code{unsigned long long}. The macros
that the installed headers read with \code{\#if} are 1 or 0, never undefined,
so that they compile under \option{-Wundef -Werror}: a program has to test
them with \code{\#if} as well, \code{\#ifdef GAOL\_PRESERVE\_ROUNDING} being
true whatever the build.

\subsection{The builds, the compilers and the options}
\label{sec:changes-builds}
\cindex{builds!differences with GAOL~4}%

The three builds agree.\newinvfive{} Before, each had its own idea (GAOL~4 had the first two, and the CMake build
is that of the 4.3 releases of GAOL v5):
\cmd{configure} optimized only when the compiler was named \cmd{g++}
exactly, computed the doubles of 32-bit x86 on the x87 unit, read the
processor of the machine building rather than the target of the compiler, and
refused macOS, Windows, s390x and riscv64; meson did not optimize, compiled
with the assertions of libstdc++, refused i686 and armv7l and linked CRlibm
by default; CMake left out \option{-msse3}, the hidden visibility and
\option{-Wconversion}. \option{-fno-fast-math}, \option{-ffp-contract=off}
and the refused compilers were in CMake only: without
\option{-ffp-contract=off}, the compilers fused the multiplications and
additions of \code{mid()}, which returned a single double, not enclosing the
midpoint of 9735 of the 10000 random intervals of a test, with meson on
macOS arm64 and with autotools on Ubuntu arm64. The three builds now write
the same macros in \file{gaol/gaol\_configuration.h}, offer the same options
(\code{GAOL\_PRESERVE\_ROUNDING}, \code{GAOL\_SIMD}, \code{GAOL\_COVERAGE},
\code{WITH\_TESTS}, \code{WITH\_EXAMPLES}, \code{BUILD\_SHARED\_LIBS} and
the build type, with their counterparts in \cmd{configure} and meson), and
decide the rest by the configuration: the fused multiply-add instructions and
the AVX-512 path by checks that run a program, the 128-bit integer by a check
that compiles. \code{GAOL\_FMA}, \code{GAOL\_PREFER\_AVX512},
\code{GAOL\_U128\_EMULATION}, \code{GAOL\_DEBUG}, \code{GAOL\_EXCEPTIONS\_ENABLED} and
\code{GAOL\_VERBOSE\_MODE} are settings of the developers, no longer options
of the user. The Debug build of each compiles GAOL with \option{-g}, without
optimization nor \code{NDEBUG}; \loption{enable-optimize} of
\cmd{configure} and \option{-Denable-optimize} of meson are gone, the Release build giving the
optimizations. The C sources of CORE-MATH are compiled with
\option{-std=gnu99}. The shared library is \file{libgaol.so.5.0.0} with the
three builds, made when asked, the static one alone by default. GAOL can be a
subproject of meson, and CMake without a build type compiles as in Release,
where, in the earlier versions of GAOL v5, \file{libgaol} had \option{-O3}
alone and CORE-MATH no optimization.
The unit tests of GAOL~4, which were in \file{check/} and needed CppUnit,
are in \file{tests/} and run without it; \cmd{make test} runs them,
\cmd{make check} the tests and the examples, \cmd{make perf} the benchmark
and \cmd{make parser} regenerates the lexer and the parser
(Chapter~\ref{chap:installation}, \pxref{sec:tests}).

\medskip
The compilers and the options that give wrong bounds are
refused.\newinvfive{} Visual C++ and clang-cl without \option{/fp:strict}
are refused, as \option{/fp:fast} is: Visual C++ then assumes rounding to
nearest, and nothing certifies the bounds. clang-cl for 32-bit x86 is given
\option{-msse2}: clang-cl 18 computes the doubles on the x87 unit there,
even with \option{/arch:SSE2}, and CORE-MATH rounds some results twice.
\option{-ffinite-math-only} is refused, as \option{-ffast-math} is; GAOL~4
refused none of them. \option{-ffast-math} and \option{-Ofast} turn it on,
and the compiler then takes NaN and infinities never to occur, in the inline
functions of the headers too: the empty interval has NaN bounds, and
$([1, 2] \cap [3, 4])$.\code{is\_empty()} is false once the compiler has
folded the test away (GCC 9.4 and 13 at \option{-O0}, \option{-O2} and \option{-O3}; Clang 18 at
\option{-O0}, and at \option{-O2} and \option{-O3} when the bounds are
\code{volatile}). With Clang, \option{-ffast-math} and \option{-Ofast} turn
\option{-frounding-math} off too, which \option{-fno-fast-math} does not turn
on again: Clang 18 compiled GAOL's operations as if rounding to nearest (54
of the 292728 checks of the test of the rounding direction failed after
\option{-ffast-math -fno-fast-math}, 126 after \option{-Ofast
-fno-fast-math}, none with all the flags after the option). What no macro
shows is not refused, and the code using GAOL is not to be compiled with it:
\option{-funsafe-math-optimizations} and \option{-fno-signed-zeros}, which
rewrite the addition by which GAOL tells the rounding direction and the
modes. The compilers that do not honour the rounding direction are refused,
the C compiler as well as the C++ compiler: Clang for 32-bit ARM processors,
and any compiler that says of \option{-frounding-math} that the rounding
mode is unsupported on its target, as Clang 14 does for 64-bit ARM. GAOL~4
accepted them; built by Clang 21 for 32-bit ARM, GAOL v5 gave 4556 of 16000
random products, squares and cubes that did not enclose their exact values
(\pxref{sec:refused}).

\medskip
GCC for Windows is given the fused multiply-add instructions (MinGW-w64,
MSYS2 and Cygwin): the three builds give it \option{-mfma
-mprefer-vector-\allowbreak width=128}.\newinvfive{} It had no flag: \option{-mfma}
turns AVX on, GCC then vectorizes with 256-bit registers and spills them to a
stack it takes to be aligned on 32 bytes, where the ABI of Windows x64
guarantees 16 (GCC bug 54412, which crashed the MinGW builds of Codac and is
fixed for GCC 17, commit 8aaa5f6e6207 of 25 September 2026), so that CORE-MATH called a software \code{fma()}. With the
second flag \cmd{objdump} finds no \code{\%ymm} in the objects, against
30598 lines with \option{-mfma} alone (GCC 16.2, MSYS2 UCRT64). On the
runner of GitHub, \code{pow(x, 3)} takes 8.7 ns rather than 52.3,
\code{log(x)} 33.6 rather than 106.4 and \code{cos(x)} 55.4 rather than
260.5; with MinGW-w64 15.2 for 32-bit x86, 12.1 ns rather than 70.5, 81.9
rather than 265.7 and 138.3 rather than 489.6. The library then needs a
processor with the FMA and AVX instructions (Intel Haswell, AMD Piledriver,
2012--2013, and later), as with the other x86 compilers. Cygwin on x86-64 is
built and tested, and refused without the instructions, the \code{fma()} of
its C library not being correctly rounded; with it, \code{previous\_float()}
and \code{next\_float()} give $-\infty$ and $+\infty$ at $\pm$\code{DBL\_MAX}
with every C library, where the \code{nextafter()} of Cygwin computed the
overflow as $x + x$ in the rounding direction of the program, so that the
lower bound of \code{sinh(-711)} did not enclose its value
(\pxref{sec:platforms}).

\medskip
The intervals of floats, \code{gaol::intervalf} and \code{gaol::interval2f},
are compiled by a developer of GAOL only.\newinvfive{} Every build compiled
and installed them, and \file{gaol/gaol} included their headers, while
neither IBEX nor Codac uses them, the manual leaves them undocumented, and
they are unfinished: \code{sqrt(intervalf)} returns its argument, and
\code{interval2f::inverse()} aborts. Their code is within
\code{\#ifdef GAOL\_FLOAT\_INTERVALS}, a macro no build has an option for.
The archive of the sources and the packages are made by CMake: \cmd{cmake
--build <build> --target package\_source} gives \file{gaol-<version>.tar.gz},
and \loption{target package} an archive of what is installed and, on Linux, a
\file{.deb}; \cmd{make dist} and \cmd{make rpm} of the autotools build are
gone, \cmd{make dist} stopping on a file it could not copy. The version is in
\file{VERSION.txt} alone, where GAOL~4 wrote it in \file{configure.ac}, in
\file{meson.build} and in two headers, and \cmd{configure --help} gives the
address for the bug reports. \code{is\_finite()} is \code{std::isfinite()}
in every build, \code{finite()} of the C library not being declared by every
one. \loption{enable-asm} and
\code{HAVE\_ROUNDING\_MATH\_OPTION} are gone, the first choosing an assembly
GAOL v5 no longer calls. CORE-MATH is in the sources, so that GAOL can be
built as a part of another project, brought in by FetchContent, with no
network access beyond its own sources. The CMake build derives from the CMake
build of GAOL in IBEX, with the compilation flags of the fork of GAOL by
Fabrice Le Bars, and the tests of \file{tests/} follow the rounding tests of
Codac (Chapter~\ref{chap:tests}). The Goldstein-Price function of the manual
and of the examples had the factor $(x + y)^2$ where the function has
$(x + y + 1)^2$, as in GAOL~4; both now compute the true function, whose
natural extension over $[-2, 2]^2$ is $[-87881320, 147125080]$, 231 times as
wide as the range, which the example says it encloses.

\section{Discussion}\label{sec:changes-discussion}
%%-------------------

The move from GAOL~4 to GAOL v5 costs a program little: the namespace to
open, the strings to read with \code{textToInterval()}, the relations
\code{==} and \code{!=} to rewrite with \code{set\_eq()},
\code{set\_disjoint()} or \code{!set\_disjoint()}, the handlers of the exceptions of GAOL~4 to
retarget on those of the standard library, and the conversions of a double or an integer
to an interval, explicit now, to write out. What the program then sees is
decided by three choices that run through the register of this chapter. The
first is CORE-MATH as the only source of the elementary functions, which gave
the tightest bounds and removed the mathematical library from the link, at
the price of sources vendored into GAOL and ported to the compilers that lack
a 128-bit integer. The second is the rounding direction taken in charge by
each operation, which made the bounds right whatever the direction the
program computes in, at the price of a check whose cost was measured on each
processor, and whose variants (the control registers written directly, the
AVX-512 path, the NEON register) each removed a part of it. The third is IEEE
1788-2015 as the arbiter wherever GAOL~4 had its own convention: the empty
set for the constructors, the comparisons, the literals read and written, the
domains of \code{log}, \code{atanh} and the roots.

\medskip
Some differences are deliberate losses: the possibly relations and
\code{certainly\_neq()}, which was one of them under another name, the timing of GAOL~4, the intervals of floats, the
headers of the sources. Others are left as they are, and the register says
so: the sign of a zero bound of a product, which depends on the kernel, the
quadratic time of the reading of a number of many digits, the inexact exception
that almost every operation raises, and the processors where neither the
flush-to-zero modes can be cleared nor \option{-mno-daz-ftz} exists. A
program that depended on one of the former behaviours finds here what
replaced it, and in the chapters this chapter refers to how the new
behaviour is defined. The appendices, after the conclusion, give where the
documentation of GAOL v5 lives and how to report what this chapter does not
cover.
